% Chapter 3, Section 6 _Linear Algebra_ Jim Hefferon
%  http://joshua.smcvt.edu/linalg.html
%  2001-Jun-12
\section{Proyeksi}
\index{proyeksi}
\noindent\textit{Bagian ini bersifat opsional.
     Bagian ini hanya menjadi prasyarat bagi
     dua bagian terakhir Bab Lima dan beberapa Topik.}

Kita telah mendeskripsikan proyeksi dari $\Re^3$ ke dalam
subruang bidang~\( xy \) sebagai pemetaan bayangan.
Gambar pertama menunjukkan alasannya,
tetapi gambar kedua juga menunjukkan bahwa sebagian bayangan jatuh ke atas.
\begin{center}  \small
  \raisebox{20.8163bp}{\includegraphics{map/mp/ch3.25}}
  \hspace*{4em}
  \raisebox{0bp}{\includegraphics{map/mp/ch3.26}}
\end{center}
\vspace*{1ex}  % graphic is poorly drawn and sticks down into text
Jadi, deskripsi yang mungkin lebih baik ialah:~proyeksi $\vec{v}$
adalah vektor $\vec{p}$ pada bidang yang memiliki sifat bahwa
seseorang yang berdiri di atas $\vec{p}$ dan memandang lurus ke atas atau
ke bawah\Dash yaitu, memandang secara ortogonal terhadap bidang\Dash
akan melihat ujung $\vec{v}$.
Dalam bagian ini kita akan memperumumnya ke proyeksi-proyeksi lain,
baik yang ortogonal maupun yang tidak ortogonal.

% The first subsection develops orthogonal projection of a vector
% into a line, which is used often in applications.
% The second subsection shows how the definition of orthogonal
% projection into a line allows us to calculate especially convenient bases
% for vector spaces, again something that often appears in applications.
% The final subsection completely generalizes projection, orthogonal or not, 
% into any subspace at all.








\subsectionoptional{Proyeksi Ortogonal ke Sebuah Garis}
Pertama-tama kita membahas proyeksi ortogonal
vektor \( \vec{v} \) ke garis $\ell$.
% We darken a point on the line if someone at that point
% looking straight up or down (that is, looking orthogonal to the line)
% will see \( \vec{v} \).
%<*PictureOrthogonalProjectionIntoALine0>
Gambar ini memperlihatkan seseorang berjalan menyusuri garis
menuju titik $\vec{p}\/$ sedemikian sehingga
ujung $\vec{v}$ tepat berada di atasnya; di sini ``di atas'' tidak berarti
sejajar dengan sumbu~$y$, melainkan ortogonal terhadap garis.
%</PictureOrthogonalProjectionIntoALine0>
\begin{center}  \small
  \includegraphics{map/mp/ch3.28}
\end{center}
%<*PictureOrthogonalProjectionIntoALine1>
Karena garis tersebut merupakan rentang suatu vektor,
$\ell=\set{c\cdot\vec{s}\suchthat c\in\Re}$,
terdapat koefisien $c_{\vec{p}}\,$
dengan sifat bahwa $\vec{v}-c_{\vec{p}}\vec{s}\,$ ortogonal terhadap
$c_{\vec{p}}\vec{s}$.
%</PictureOrthogonalProjectionIntoALine1>
\begin{center}  \small
  \includegraphics{map/mp/ch3.29}
\end{center}
%<*PictureOrthogonalProjectionIntoALine2>
Untuk menentukan koefisien ini, perhatikan bahwa karena
$\vec{v}-c_{\vec{p}}\vec{s}\,$ ortogonal terhadap suatu kelipatan skalar
dari $\vec{s}$, vektor itu pasti ortogonal terhadap $\vec{s}\,$ sendiri.
Dengan demikian,
$(\vec{v}-c_{\vec{p}}\vec{s})\dotprod \vec{s}=0$ memberikan
$c_{\vec{p}}=\vec{v}\dotprod \vec{s}/\vec{s}\dotprod \vec{s}$.
%</PictureOrthogonalProjectionIntoALine2>

\begin{definition}  \label{def:ProjIntoLine}
%<*df:ProjIntoLine>
\definend{Proyeksi ortogonal \( \vec{v} \) ke garis yang direntang oleh
\( \vec{s}\, \) tak nol}\index{proyeksi!ortogonal}%
\index{proyeksi!ke sebuah garis} adalah vektor berikut.
\begin{equation*}
  \proj{\vec{v}}{\spanof{\vec{s}\,}}=
  \frac{ \vec{v}\dotprod\vec{s} }{ \vec{s}\dotprod\vec{s} }\cdot\vec{s}
\end{equation*}
%</df:ProjIntoLine>
\end{definition}

% \noindent \nearbyexercise{exer:LineProjFormIndOfVec} checks that the 
% outcome of the calculation depends only on the line and not on which vector 
% \( \vec{s}\, \) we used to describe that line.
\noindent (Di sini digunakan ungkapan `direntang oleh \( \vec{s}\, \)'
alih-alih ungkapan yang lebih formal `rentang himpunan \( \set{\vec{s}\,} \)'.
Ungkapan yang lebih santai ini lazim digunakan.)

\begin{example}
Untuk memproyeksikan secara ortogonal
vektor $\binom{2}{3}$ ke garis \( y=2x \),
pertama-tama pilih suatu vektor arah bagi garis tersebut.
\begin{equation*}
  \vec{s}=\colvec[r]{1 \\ 2}
\end{equation*}
Perhitungannya mudah.
\begin{center}  \small
  \begin{tabular}{@{}c@{}}\includegraphics{map/mp/ch3.30}\end{tabular}
  \qquad
  $\frac{ \colvec[r]{2 \\ 3}\dotprod\colvec[r]{1 \\ 2} }{% 
          \colvec[r]{1 \\ 2}\dotprod\colvec[r]{1 \\ 2} }
     \cdot\colvec[r]{1 \\ 2}
    ={\displaystyle \frac{8}{5}}\cdot\colvec[r]{1 \\ 2}=\colvec[r]{8/5 \\ 16/5}$
\end{center}
\end{example}

\begin{example}
Dalam \( \Re^3 \), proyeksi ortogonal vektor umum
\begin{equation*}
  \colvec{x \\ y \\ z}
\end{equation*}
ke sumbu~\( y \) adalah
\begin{equation*}
  \frac{ \colvec{x \\ y \\ z}\dotprod\colvec[r]{0 \\ 1 \\ 0} }{
         \colvec[r]{0 \\ 1 \\ 0}\dotprod\colvec[r]{0 \\ 1 \\ 0} }
          \cdot\colvec[r]{0 \\ 1 \\ 0}=\colvec{0 \\ y \\ 0}
\end{equation*}
yang sesuai dengan dugaan intuitif kita.
\end{example}

Gambar di atas, yang memperlihatkan seseorang berjalan menyusuri garis hingga
ujung $\vec{v}$ berada tepat di atasnya, merupakan salah satu cara memandang
proyeksi ortogonal suatu vektor ke sebuah garis.
Kita menutup subbagian ini dengan dua cara pandang lain.


\begin{example} \label{exam:RailCar}
Sebuah gerbong kereta yang dibiarkan tanpa rem pada jalur timur--barat didorong
angin yang bertiup ke timur laut dengan kecepatan lima belas mil per jam;
berapakah kecepatan yang akan dicapai gerbong tersebut?
\begin{center}  \small
  \includegraphics{map/mp/ch3.31}
\end{center}
Untuk angin tersebut kita gunakan vektor dengan panjang $15$ yang menunjuk ke
timur laut.
\begin{equation*}
  \vec{v}=\colvec[r]{15\sqrt{1/2} \\ 15\sqrt{1/2}}
\end{equation*}
Gerbong hanya dipengaruhi oleh bagian angin yang bertiup dalam arah
timur--barat\Dash bagian $\vec{v}$ dalam arah sumbu~$x$ adalah berikut ini
(gambar ini menggunakan sudut pandang yang sama dengan gambar gerbong di atas).
\begin{center}
  \begin{tabular}{@{}c@{}}\includegraphics{map/mp/ch3.32}\end{tabular}
  \qquad
  $\displaystyle \vec{p}=\colvec{15\sqrt{1/2} \\ 0}$
\end{center}
Jadi, gerbong akan mencapai kecepatan
\( 15\sqrt{1/2} \) mil per jam ke arah timur.
\end{example}

Jadi, cara lain memandang
gambar sebelum definisi adalah bahwa gambar itu memperlihatkan
$\vec{v}$ yang diuraikan menjadi dua bagian: bagian $\vec{p}$ yang searah
dengan garis, dan bagian yang ortogonal terhadap garis
(ditunjukkan di atas pada sumbu utara--selatan).
Kedua bagian ini tidak saling berinteraksi,
dalam arti bahwa gerbong timur--barat sama sekali tidak dipengaruhi oleh
bagian angin utara--selatan (lihat \nearbyexercise{exer:PerpAreInd}).
Jadi, proyeksi ortogonal \( \vec{v} \) ke garis yang direntang oleh
\( \vec{s}\, \) dapat dipandang sebagai bagian dari \( \vec{v}\, \)
yang terletak dalam arah \( \vec{s} \).

Cara berguna lain untuk memandang proyeksi ortogonal ke sebuah garis
adalah dengan meminta orang tersebut berdiri di atas vektor, bukan di atas garis.
Orang ini memegang tali yang dilingkarkan mengitari garis.
Ketika ia menariknya, lingkaran tali bergeser pada garis.
\begin{center}  \small
  \includegraphics{map/mp/ch3.33}
\end{center}
Ketika tali menegang, tali itu ortogonal terhadap garis.
Artinya, proyeksi \( \vec{p}\, \) dapat dipandang sebagai vektor pada garis
yang paling dekat dengan \( \vec{v} \)
(lihat \nearbyexercise{exer:ProjMinimizesDist}).

\begin{example}
Sebuah kapal selam melacak kapal yang bergerak sepanjang garis \( y=3x+2 \).
Jangkauan torpedonya setengah mil.
Jika kapal selam tetap berada di tempatnya, yakni di titik asal pada peta di
bawah, apakah kapal tersebut akan melintas dalam jangkauan?
\begin{center}  \small
  \includegraphics{map/mp/ch3.34}
\end{center} 
Rumus proyeksi
ke sebuah garis tidak langsung berlaku karena garis itu tidak melalui
titik asal, sehingga bukan merupakan rentang dari suatu $\vec{s}$.
Untuk menyesuaikannya, mula-mula kita geser seluruh peta dua satuan ke bawah.
Kini garisnya adalah $y=3x$, suatu subruang.
Kita melakukan proyeksi untuk memperoleh
titik $\vec{p}$ pada garis yang paling dekat dengan
\begin{equation*}
  \vec{v}=\colvec[r]{0 \\ -2}
\end{equation*}
posisi kapal selam setelah digeser.
\begin{equation*}
  \vec{p}=\frac{\colvec[r]{0 \\ -2}\dotprod\colvec[r]{1 \\ 3}}{
                \colvec[r]{1 \\ 3}\dotprod\colvec[r]{1 \\ 3}}
          \cdot \colvec[r]{1 \\ 3}=\colvec[r]{-3/5 \\ -9/5}
\end{equation*}
Jarak antara $\vec{v}$ dan $\vec{p}$ kira-kira
\( 0.63 \)~mil.
Kapal tersebut tidak akan pernah berada dalam jangkauan.
\end{example}


\begin{exercises}
  \recommended \item
    Proyeksikan vektor pertama secara ortogonal
    ke garis yang direntang oleh vektor kedua.
    \begin{exparts*}
      \partsitem \mbox{\( \colvec[r]{2 \\ 1} \), \( \colvec[r]{3 \\ -2} \)}
      \partsitem \mbox{\( \colvec[r]{2 \\ 1} \), \( 
            \colvec[r]{3 \\ 0} \)}
      \partsitem \mbox{\( \colvec[r]{1 \\ 1 \\ 4} \), \( 
         \colvec[r]{1 \\ 2 \\ -1} \)}
      \partsitem \mbox{\( \colvec[r]{1 \\ 1 \\ 4} \), \( 
         \colvec[r]{3 \\ 3 \\ 12} \)}
    \end{exparts*}
    \begin{answer}
       % Each is a straightforward application of the formula from 
       % \nearbydefinition{def:ProjIntoLine}.
       \begin{exparts}
        \partsitem 
          $\displaystyle \frac{\colvec[r]{2 \\ 1}\dotprod\colvec[r]{3 \\ -2}}{%
                               \colvec[r]{3 \\ -2}\dotprod\colvec[r]{3 \\ -2}}
             \cdot\colvec[r]{3 \\ -2}
           =\frac{4}{13}
             \cdot\colvec[r]{3 \\ -2}
           =\colvec[r]{12/13 \\ -8/13}$
        \partsitem 
          $\displaystyle \frac{\colvec[r]{2 \\ 1}\dotprod\colvec[r]{3 \\ 0}}{%
                               \colvec[r]{3 \\ 0}\dotprod\colvec[r]{3 \\ 0}}
             \cdot\colvec[r]{3 \\ 0}
           =\frac{2}{3}
             \cdot\colvec[r]{3 \\ 0}
           =\colvec[r]{2 \\ 0}$
        \partsitem 
          $\displaystyle 
             \frac{\colvec[r]{1 \\ 1 \\ 4}\dotprod\colvec[r]{1 \\ 2 \\ -1}}{%
                   \colvec[r]{1 \\ 2 \\ -1}\dotprod\colvec[r]{1 \\ 2 \\ -1}}
             \cdot\colvec[r]{1 \\ 2 \\ -1}
           =\frac{-1}{6}
             \cdot\colvec[r]{1 \\ 2 \\ -1}
           =\colvec[r]{-1/6 \\ -1/3 \\ 1/6}$
        \partsitem 
          $\displaystyle 
             \frac{\colvec[r]{1 \\ 1 \\ 4}\dotprod\colvec[r]{3 \\ 3 \\ 12}}{%
                   \colvec[r]{3 \\ 3 \\ 12}\dotprod\colvec[r]{3 \\ 3 \\ 12}}
             \cdot\colvec[r]{3 \\ 3 \\ 12}
           =\frac{1}{3}
             \cdot\colvec[r]{3 \\ 3 \\ 12}
           =\colvec[r]{1 \\ 1 \\ 4}$
      \end{exparts}  
    \end{answer}
  \recommended \item 
    Proyeksikan vektor tersebut secara ortogonal ke garis.
    \begin{exparts*}
      \partsitem 
       $\colvec[r]{2 \\ -1 \\ 4},~\set{c\colvec[r]{-3 \\ 1 \\ -3}\suchthat c\in\Re}$
      \partsitem $\colvec[r]{-1 \\ -1}$,~garis $y=3x$
    \end{exparts*}
    \begin{answer}
      \begin{exparts}
        \partsitem 
          $\displaystyle
             \frac{\colvec[r]{2 \\ -1 \\ 4}\dotprod\colvec[r]{-3 \\ 1 \\ -3}}{%
                   \colvec[r]{-3 \\ 1 \\ -3}\dotprod\colvec[r]{-3 \\ 1 \\ -3}}
             \cdot\colvec[r]{-3 \\ 1 \\ -3}
           =\frac{-19}{19}\cdot\colvec[r]{-3 \\ 1 \\ -3}
           =\colvec[r]{3 \\ -1 \\ 3}$
         \partsitem Dengan menuliskan garis sebagai
           \begin{equation*} 
             \set{c\cdot \colvec[r]{1 \\ 3}\suchthat c\in\Re}
           \end{equation*}
           diperoleh proyeksi berikut.
           \begin{equation*}
             \frac{\colvec[r]{-1 \\ -1}\dotprod\colvec[r]{1 \\ 3}}{%
                   \colvec[r]{1 \\ 3}\dotprod\colvec[r]{1 \\ 3}}
              \cdot\colvec[r]{1 \\ 3}
             =\frac{-4}{10}\cdot\colvec[r]{1 \\ 3}
             =\colvec[r]{-2/5 \\ -6/5}
           \end{equation*}
      \end{exparts}
    \end{answer}
  \item 
    Meskipun gambar memandu pengembangan
    \nearbydefinition{def:ProjIntoLine}, kita tidak terbatas pada ruang yang
    dapat kita gambar.
    Dalam $\Re^4$, proyeksikan vektor ini ke garis berikut.
    \begin{equation*}
      \vec{v}=\colvec[r]{1 \\ 2 \\ 1 \\ 3}
      \qquad
      \ell=\set{c\cdot\colvec[r]{-1 \\ 1 \\ -1 \\ 1}\suchthat c\in\Re}
    \end{equation*}
    \begin{answer}
      $\displaystyle 
         \frac{\colvec[r]{1 \\ 2 \\ 1 \\ 3}\dotprod\colvec[r]{-1 \\ 1 \\ -1 \\ 1}}{%
               \colvec[r]{-1 \\ 1 \\ -1 \\ 1}\dotprod\colvec[r]{-1 \\ 1 \\ -1 \\ 1}}
         \cdot\colvec[r]{-1 \\ 1 \\ -1 \\ 1}
       =\frac{3}{4}
         \cdot\colvec[r]{-1 \\ 1 \\ -1 \\ 1}
       =\colvec[r]{-3/4 \\ 3/4 \\ -3/4 \\ 3/4}$ 
    \end{answer}
  \recommended \item 
    \nearbydefinition{def:ProjIntoLine} menggunakan dua vektor, $\vec{s}$ dan
    $\vec{v}$.
    Tinjau transformasi pada $\Re^2$ yang diperoleh dengan menetapkan
    \begin{equation*}
      \vec{s}=\colvec[r]{3 \\ 1}
    \end{equation*}
    lalu memproyeksikan $\vec{v}$ ke garis yang merupakan rentang $\vec{s}$.
    Terapkan transformasi itu pada vektor-vektor berikut.
    \begin{exparts*}
      \partsitem $\colvec[r]{1 \\ 2}$
      \partsitem $\colvec[r]{0 \\ 4}$ 
    \end{exparts*}
    Tunjukkan bahwa secara umum transformasi proyeksinya adalah berikut.
    \begin{equation*}
      \colvec{x_1 \\ x_2}
      \mapsto
      \colvec{(9x_1+3x_2)/10 \\ (3x_1+x_2)/10}
    \end{equation*}
    Nyatakan tindakan transformasi ini dengan sebuah matriks.
    \begin{answer}
      \begin{exparts}
        \partsitem $\displaystyle
          \frac{\colvec[r]{1 \\ 2}\dotprod\colvec[r]{3 \\ 1}}{%
                 \colvec[r]{3 \\ 1}\dotprod\colvec[r]{3 \\ 1}}
          \cdot \colvec[r]{3 \\ 1}
          =\frac{1}{2}\cdot\colvec[r]{3 \\ 1}
          =\colvec[r]{3/2 \\ 1/2}$
        \partsitem $\displaystyle
          \frac{\colvec[r]{0 \\ 4}\dotprod\colvec[r]{3 \\ 1}}{%
                 \colvec[r]{3 \\ 1}\dotprod\colvec[r]{3 \\ 1}}
          \cdot \colvec[r]{3 \\ 1}
          =\frac{2}{5}\cdot\colvec[r]{3 \\ 1}
          =\colvec[r]{6/5 \\ 2/5}$
      \end{exparts}
      \noindent Secara umum proyeksinya adalah berikut.
      \begin{equation*}
          \frac{\colvec{x_1 \\ x_2}\dotprod\colvec[r]{3 \\ 1}}{%
                 \colvec[r]{3 \\ 1}\dotprod\colvec[r]{3 \\ 1}}
          \cdot \colvec[r]{3 \\ 1}
          =\frac{3x_1+x_2}{10}\cdot\colvec[r]{3 \\ 1}
          =\colvec{(9x_1+3x_2)/10 \\ (3x_1+x_2)/10}
      \end{equation*}
      Matriks yang sesuai adalah berikut.
      \begin{equation*}
        \begin{mat}[r]
          9/10  &3/10  \\
          3/10  &1/10
        \end{mat}
      \end{equation*}  
     \end{answer}
  \item \label{exer:PerpAreInd}
    \nearbyexample{exam:RailCar} menunjukkan bahwa proyeksi menguraikan
    $\vec{v}$ menjadi dua bagian, $\proj{\vec{v}\,}{\spanof{\vec{s}\,}}$ dan
    $\vec{v}-\proj{\vec{v}\,}{\spanof{\vec{s}\,}}$, yang tidak saling
    berinteraksi.
    Ingat bahwa keduanya ortogonal.
    Tunjukkan bahwa sebarang dua vektor ortogonal tak nol membentuk himpunan
    yang bebas linear.
    \begin{answer}
      Misalkan $\vec{v}_1$ dan $\vec{v}_2$ tak nol dan ortogonal.
      Tinjau hubungan linear
      $c_1\vec{v}_1+c_2\vec{v}_2=\zero$.
      Ambil hasil kali titik kedua ruas persamaan dengan $\vec{v}_1$ untuk
      memperoleh
      \begin{multline*}
        \vec{v}_1\dotprod(c_1\vec{v}_1+c_2\vec{v}_2)
        =c_1\cdot(\vec{v}_1\dotprod\vec{v}_1)
           +c_2\cdot(\vec{v}_1\dotprod\vec{v}_2)            \\
        =c_1\cdot (\vec{v}_1\dotprod\vec{v}_1)+c_2\cdot 0
        =c_1\cdot (\vec{v}_1\dotprod\vec{v}_1)
      \end{multline*}
      yang sama dengan $\vec{v}_1\dotprod\zero=\zero$.
      Karena diasumsikan bahwa $\vec{v}_1$ tak nol, hal ini memberikan bahwa
      $c_1$ nol.
      Pembuktian bahwa $c_2$ nol serupa.
    \end{answer}
  \item    
    \begin{exparts}
      \partsitem Apa proyeksi ortogonal \( \vec{v} \) ke sebuah garis jika
        \( \vec{v} \) merupakan anggota garis tersebut?
      \partsitem Tunjukkan bahwa jika $\vec{v}$ bukan anggota garis tersebut,
        maka himpunan
        $\set{\vec{v},\vec{v}-\proj{\vec{v}\,}{\spanof{\vec{s}\,}}}$
        bebas linear.
    \end{exparts}
    \begin{answer}
      \begin{exparts}
       \partsitem Jika vektor $\vec{v}\,$ berada pada garis tersebut, maka
         proyeksi ortogonalnya adalah \( \vec{v} \).
         Untuk memverifikasinya melalui perhitungan, perhatikan bahwa karena
         $\vec{v}\,$ berada pada garis tersebut, berlaku
         \( \vec{v}=c_{\vec{v}}\cdot\vec{s}\, \) bagi suatu skalar
         $c_{\vec{v}}$.
         \begin{equation*}
           \frac{\vec{v}\dotprod\vec{s}}{\vec{s}\dotprod\vec{s}}\cdot\vec{s}
           =\frac{c_{\vec{v}}\cdot \vec{s}\dotprod\vec{s}}{%
                         \vec{s}\dotprod\vec{s}}\cdot\vec{s}
           =c_{\vec{v}}\cdot\frac{\vec{s}\dotprod\vec{s}}{%
                                  \vec{s}\dotprod\vec{s}}\cdot\vec{s}
           =c_{\vec{v}}\cdot 1 \cdot\vec{s}
           =\vec{v}
         \end{equation*}
          (\textit{Catatan}.
          Jika kita mengasumsikan bahwa $\vec{v}\,$ tak nol, perhitungan di
          atas dapat disederhanakan dengan mengambil $\vec{s}\,$ sebagai
          $\vec{v}$.)
       \partsitem Tuliskan $c_{\vec{p}}\vec{s}\,$ untuk proyeksi
         $\proj{\vec{v}}{\spanof{\vec{s}\,}}$. 
         Perhatikan bahwa, berdasarkan asumsi bahwa $\vec{v}$ tidak berada
         pada garis, baik $\vec{v}$ maupun
         $\vec{v}-c_{\vec{p}}\vec{s}\,$ tak nol.
         Perhatikan pula bahwa jika $c_{\vec{p}}\,$ nol, sesungguhnya kita
         sedang meninjau himpunan beranggota tunggal $\set{\vec{v}\,}$;
         karena $\vec{v}$ tak nol, himpunan ini pasti bebas linear.
         Oleh karena itu, tinggal ditinjau kasus bahwa $c_{\vec{p}}$ tak nol.

         Dengan menyusun hubungan linear
         \begin{equation*}
           a_1(\vec{v})+a_2(\vec{v}-c_{\vec{p}}\vec{s})=\zero
         \end{equation*}
         diperoleh persamaan
         $(a_1+a_2)\cdot\vec{v}=a_2c_{\vec{p}}\cdot\vec{s}$.
         Karena $\vec{v}\,$ tidak berada pada garis, kedua skalar
         $a_1+a_2$ dan $a_2 c_{\vec{p}}$ harus nol.
         Kasus $c_{\vec{p}}=0$ telah ditangani di atas, sehingga pada kasus
         yang tersisa $a_2=0$, dan ini juga memberikan $a_1=0$.
         Jadi, himpunan tersebut bebas linear.
     \end{exparts}
    \end{answer}
  \item 
    \nearbydefinition{def:ProjIntoLine} mensyaratkan bahwa $\vec{s}\,$ tak nol.
    Mengapa?
    Apa definisi yang tepat bagi proyeksi ortogonal suatu vektor ke garis
    (degenerat) yang direntang oleh vektor nol?
    \begin{answer}
      Jika $\vec{s}\,$ adalah vektor nol, maka ungkapan
          \begin{equation*}
            \proj{\vec{v}}{\spanof{\vec{s}\,}}=
            \frac{ \vec{v}\dotprod\vec{s} }{%
                    \vec{s}\dotprod\vec{s} }\cdot\vec{s}
          \end{equation*}
      memuat pembagian dengan nol sehingga tidak terdefinisi.
      Adapun definisi yang tepat: agar proyeksi berada dalam rentang vektor
      nol, proyeksi itu harus didefinisikan sebagai \( \zero \).
    \end{answer}
  \item 
    Apakah setiap vektor merupakan proyeksi suatu vektor lain ke suatu garis?
    \begin{answer}
      Setiap vektor dalam \( \Re^n \) merupakan proyeksi suatu vektor lain ke
      sebuah garis, asalkan dimensi \( n \) lebih besar dari satu.
      (Jelas bahwa setiap vektor merupakan proyeksi dirinya sendiri ke garis
      yang memuatnya; pertanyaannya ialah mencari suatu vektor selain
      $\vec{v}$ yang proyeksinya adalah \( \vec{v} \).)

      Misalkan \( \vec{v}\in\Re^n \) dengan \( n>1 \).
      Jika \( \vec{v}=\zero \), ambil garis
      \( \ell=\spanof{\vec{e}_1} \) dan vektor
      \( \vec{w}=\vec{e}_2 \). Karena \( n>1 \), vektor ini tersedia,
      berbeda dari \( \vec{v} \), dan proyeksinya ke \( \ell \) adalah
      \( \zero=\vec{v} \).
      Jadi, untuk bagian pembuktian yang tersisa, misalkan
      \( \vec{v}\neq\zero \) dan ambil garis
      \( \ell=\set{c\vec{v}\suchthat c\in\Re} \).
      Misalkan \( v_1,\dots,v_n \) adalah komponen-komponen \( \vec{v} \);
      karena \( n>1 \), setidaknya ada dua komponen.
      Jika suatu \( v_i \) nol, maka vektor \( \vec{w}=\vec{e}_i \) tegak
      lurus terhadap \( \vec{v} \).
      Jika tidak ada komponennya yang nol, maka vektor \( \vec{w}\, \) yang
      komponennya \( v_2,-v_1,0,\dots ,0 \) tegak lurus terhadap
      \( \vec{v} \).
      Perhatikan bahwa \( \vec{v}+\vec{w} \) tidak
      sama dengan \( \vec{v} \), dan bahwa \( \vec{v} \) merupakan proyeksi
      \( \vec{v}+\vec{w} \) ke \( \ell \).
      \begin{equation*}
        \frac{ (\vec{v}+\vec{w})\dotprod\vec{v} }{ 
               \vec{v}\dotprod\vec{v} }\cdot\vec{v}
        =\bigl(\frac{ \vec{v}\dotprod\vec{v} }{ 
                     \vec{v}\dotprod\vec{v} }+
               \frac{ \vec{w}\dotprod\vec{v} }{ \
                      \vec{v}\dotprod\vec{v} }\bigr)\cdot\vec{v}
        =\frac{ \vec{v}\dotprod\vec{v} }{ \vec{v}\dotprod\vec{v} }\cdot\vec{v}
        =\vec{v}
      \end{equation*}

      Tinggal kita tangani kasus $n=0$ dan $n=1$.
      Kasus berdimensi \( n=0 \) adalah ruang vektor trivial; di sini hanya
      ada satu vektor sehingga vektor itu tidak dapat dinyatakan sebagai
      proyeksi vektor yang berbeda.
      Dalam kasus berdimensi $n=1$ hanya ada satu garis (tak degenerat), dan
      setiap vektor berada di dalamnya; jadi setiap vektor hanya merupakan
      proyeksi dirinya sendiri.
    \end{answer}
  \item 
    Tunjukkan bahwa panjang proyeksi \( \vec{v}\, \) ke garis yang direntang
    oleh \( \vec{s}\, \) sama dengan nilai mutlak bilangan
    \( \vec{v}\dotprod\vec{s}\, \) dibagi panjang vektor \( \vec{s}\, \).
    \begin{answer}
      Buktinya berupa suatu perhitungan.
      Ingat bahwa bagi sebarang vektor $\vec{u}$, panjangnya ditentukan oleh
      $\absval{\vec{u}}^2=\vec{u}\dotprod\vec{u}$.
      Jadi, berikut adalah kuadrat panjang proyeksi tersebut.
      \begin{equation*}
        \absval{\frac{\vec{v}\dotprod\vec{s}}{\vec{s}\dotprod\vec{s}}
              \cdot\vec{s}\,}^2
        =
        \bigl(\frac{\vec{v}\dotprod\vec{s}}{\vec{s}\dotprod\vec{s}}
           \cdot \vec{s}\bigr)
        \dotprod
        \bigl(\frac{\vec{v}\dotprod\vec{s}}{\vec{s}\dotprod\vec{s}}
           \cdot \vec{s}\bigr)
        =
        \bigl(\frac{\vec{v}\dotprod\vec{s}}{\vec{s}\dotprod\vec{s}}\bigr)^2
        \cdot(\vec{s}\dotprod\vec{s})
        =
        \frac{(\vec{v}\dotprod\vec{s})^2}{(\vec{s}\dotprod\vec{s})^2}
        \cdot(\vec{s}\dotprod\vec{s})
        =
        \frac{(\vec{v}\dotprod\vec{s})^2}{\absval{\vec{s}}^2}
      \end{equation*} 
    \end{answer}
  \item
    Temukan rumus jarak dari sebuah titik ke sebuah garis.
    \begin{answer}
      Karena proyeksi \( \vec{v} \) ke garis yang direntang oleh
      \( \vec{s} \) adalah
      \begin{equation*}
        \frac{\vec{v}\dotprod\vec{s}}{\vec{s}\dotprod\vec{s}}\cdot\vec{s}
      \end{equation*}
      kuadrat jarak dari titik ke garis adalah sebagai berikut.
      \begin{align*}
        \absval{\vec{v}-
                \frac{\vec{v}\dotprod\vec{s}}{\vec{s}\dotprod\vec{s}}
                \cdot\vec{s}\, }^2
        &=\big( \vec{v}-
                \frac{\vec{v}\dotprod\vec{s}}{\vec{s}\dotprod\vec{s}}
                \cdot\vec{s} \big)
          \dotprod
          \big( \vec{v}-
                \frac{\vec{v}\dotprod\vec{s}}{\vec{s}\dotprod\vec{s}}
                \cdot\vec{s} \big)                                  \\
        &=\vec{v}\dotprod\vec{v}
        -\vec{v}\dotprod(
              \frac{\vec{v}\dotprod\vec{s}}{\vec{s}\dotprod\vec{s}}\cdot\vec{s}
           )
        -(
              \frac{\vec{v}\dotprod\vec{s}}{\vec{s}\dotprod\vec{s}}
               \cdot\vec{s}\,
         )\dotprod\vec{v}
        +\big(\frac{\vec{v}\dotprod\vec{s}}{
                \vec{s}\dotprod\vec{s}}\cdot\vec{s}\,\big)
          \dotprod
          \big(\frac{\vec{v}\dotprod\vec{s}}{
                \vec{s}\dotprod\vec{s}}\cdot\vec{s}\,\big)     \\   
        &=\vec{v}\dotprod\vec{v}
        -2\cdot\big(\frac{\vec{v}\dotprod\vec{s}}{\vec{s}\dotprod\vec{s}}\big)
           \cdot\vec{v}\dotprod\vec{s}
        +\big(\frac{\vec{v}\dotprod\vec{s}}{
                \vec{s}\dotprod\vec{s}}\big)^2\cdot(\vec{s}\dotprod\vec{s})  \\
        &=\frac{(\vec{v}\dotprod\vec{v}\,)\cdot(\vec{s}\dotprod\vec{s}\,)
                 -2\cdot(\vec{v}\dotprod\vec{s}\,)^2
                 +(\vec{v}\dotprod\vec{s}\,)^2}{%
               \vec{s}\dotprod\vec{s}}                               \\     
        &=\frac{(\vec{v}\dotprod\vec{v}\,)(\vec{s}\dotprod\vec{s}\,)
                -(\vec{v}\dotprod\vec{s}\,)^2}{\vec{s}\dotprod\vec{s}}
      \end{align*}  
    \end{answer}
  \item 
    Dengan menggunakan Kalkulus, temukan skalar \( c \) sedemikian sehingga
    titik \( (cs_1,cs_2) \) berjarak minimum dari titik \( (v_1,v_2) \)
    (yakni, tinjau fungsi jarak, samakan turunan pertamanya dengan nol, lalu
    selesaikan).
    Perumum hasilnya ke \( \Re^n \).
    \begin{answer}
      Karena akar kuadrat merupakan fungsi yang naik secara ketat, kita dapat
      meminimumkan \( d(c)=(cs_1-v_1)^2+(cs_2-v_2)^2 \), bukan akar kuadrat
      dari \( d \).
      Turunannya adalah
      $dd/dc=2(cs_1-v_1)\cdot s_1 +2(cs_2-v_2)\cdot s_2$.
      Menyamakannya dengan nol,
      $2(cs_1-v_1)\cdot s_1 +2(cs_2-v_2)\cdot s_2
        =c\cdot(2s_1^2+2s_2^2)-2(v_1s_1+v_2s_2)=0$
      memberikan satu-satunya titik kritis.
      \begin{equation*}
        c=\frac{v_1s_1+v_2s_2}{{s_1}^2+{s_2}^2}
         =\frac{\vec{v}\dotprod\vec{s}}{\vec{s}\dotprod\vec{s}}
      \end{equation*}
      Sekarang turunan kedua terhadap $c$,
      \begin{equation*}
        \frac{d^2\,d}{dc^2}=2{s_1}^2+2{s_2}^2
      \end{equation*}
      positif secara ketat (selama \( s_1 \) dan \( s_2 \) tidak keduanya nol;
      jika keduanya nol, soalnya trivial), sehingga titik kritis tersebut
      merupakan suatu minimum.

      Perumuman ke $\Re^n$ langsung.
      Tinjau $d_n(c)=(cs_1-v_1)^2+\dots+(cs_n-v_n)^2$, ambil turunannya, dan
      seterusnya.
    \end{answer}
  \recommended \item \label{exer:ProjMinimizesDist}
    Misalkan $\vec{p}$ adalah proyeksi ortogonal
    $\vec{v}\in\Re^n$ ke garis~$\ell$.
    Tunjukkan bahwa $\vec{p}$ adalah titik pada garis yang paling dekat dengan
    $\vec{v}$.
    \begin{answer}
      Misalkan $\vec{w}\in\ell$.
      Karena ini merupakan proyeksi ortogonal, kedua vektor
      $\vec{v}-\vec{p}$ dan $\vec{p}-\vec{w}$ saling tegak lurus.
      Teorema Pythagoras berlaku, sehingga sisi miring
      $\vec{v}-\vec{w}$ setidaknya sepanjang $\vec{v}-\vec{p}$.
    \end{answer}
  \item
    Buktikan bahwa panjang proyeksi ortogonal suatu vektor ke sebuah garis
    kurang dari atau sama dengan panjang vektor tersebut.
    \begin{answer}
      Bagi sebarang vektor $\vec{u}$, kuadrat panjangnya diberikan oleh
      $\absval{\vec{u}}^2=\vec{u}\dotprod\vec{u}$.
      Jadi, berikut adalah kuadrat panjang proyeksinya.
      \begin{equation*}
        \absval{\frac{\vec{v}\dotprod\vec{s}}{\vec{s}\dotprod\vec{s}}
              \cdot\vec{s}\,}^2
        =
        \bigl(\frac{\vec{v}\dotprod\vec{s}}{\vec{s}\dotprod\vec{s}}
           \cdot \vec{s}\,\bigr)
        \dotprod
        \bigl(\frac{\vec{v}\dotprod\vec{s}}{\vec{s}\dotprod\vec{s}}
           \cdot \vec{s}\,\bigr)
        =
        \bigl(\frac{\vec{v}\dotprod\vec{s}}{\vec{s}\dotprod\vec{s}}\bigr)^2
        \cdot(\vec{s}\dotprod\vec{s})
        =
        \frac{(\vec{v}\dotprod\vec{s})^2}{(\vec{s}\dotprod\vec{s})^2}
        \cdot(\vec{s}\dotprod\vec{s})
        =
        \frac{(\vec{v}\dotprod\vec{s})^2}{\absval{\vec{s}\,}^2}
      \end{equation*} 
      Dengan demikian, panjangnya adalah
      $\absval{\vec{v}\dotprod\vec{s}}/\absval{\vec{s}\,}$, yakni nilai mutlak
      $\vec{v}\dotprod\vec{s}$ dibagi panjang~$\vec{s}$.
      Ketaksamaan Cauchy--Schwarz,
      $\absval{\vec{v}\dotprod\vec{s}\,}
         \leq\absval{\vec{v}\,}\cdot\absval{\vec{s}\,}$, 
      memberikan
      $\absval{\vec{v}\dotprod\vec{s}}/\absval{\vec{s}\,}\leq
        \absval{\vec{v}\,}\cdot\absval{\vec{s}\,}/\absval{\vec{s}\,}
        =\absval{\vec{v}\,}$.
    \end{answer}
  \recommended \item  \label{exer:LineProjFormIndOfVec}
    Tunjukkan bahwa definisi proyeksi ortogonal ke sebuah garis tidak
    bergantung pada vektor perentangnya:~jika \( \vec{s} \) merupakan kelipatan
    tak nol dari \( \vec{q} \), maka
    \( (\vec{v}\dotprod\vec{s}/\vec{s}\dotprod\vec{s}\,)\cdot\vec{s} \) sama dengan
    \( (\vec{v}\dotprod\vec{q}/\vec{q}\dotprod\vec{q}\,)\cdot\vec{q} \).
    \begin{answer}
       Tuliskan \( c\vec{s} \) sebagai \( \vec{q} \), lalu hitung:
       \( (\vec{v}\dotprod c\vec{s}/c\vec{s}\dotprod c\vec{s}\,)\cdot c\vec{s}=
          (\vec{v}\dotprod \vec{s}/\vec{s}\dotprod \vec{s}\,)\cdot \vec{s} \).
    \end{answer}
  \item
    Tinjau fungsi dari bidang ke dirinya sendiri yang membawa suatu vektor ke
    proyeksinya pada garis \( y=x \).
    Dua bagian berikut masing-masing menunjukkan bahwa pemetaan ini linear:
    yang pertama dengan cara yang terikat koordinat (yakni, menetapkan sebuah
    basis lalu menghitung), dan yang kedua dengan cara yang lebih konseptual.
    \begin{exparts} 
      \partsitem Berikan matriks yang mendeskripsikan tindakan fungsi tersebut.
      \partsitem Tunjukkan bahwa pemetaan ini dapat diperoleh dengan mula-mula
        memutar seluruh bidang sebesar \( \pi/4 \)~radian searah jarum jam,
        kemudian memproyeksikannya ke sumbu~\( x \), lalu memutarnya kembali
        sebesar \( \pi/4 \)~radian berlawanan arah jarum jam.
    \end{exparts}
    \begin{answer}
      \begin{exparts}
        \partsitem Dengan menetapkan
          \begin{equation*}
            \vec{s}=\colvec[r]{1 \\ 1}
          \end{equation*}
          sebagai vektor yang rentangnya adalah garis tersebut, rumusnya
          memberikan tindakan berikut,
          \begin{equation*}
            \colvec{x \\ y}
             \mapsto
            \frac{\colvec{x \\ y}\dotprod\colvec[r]{1 \\ 1}}{%
                  \colvec[r]{1 \\ 1}\dotprod\colvec[r]{1 \\ 1}}\cdot\colvec[r]{1 \\ 1}
            =\frac{x+y}{2}\cdot\colvec[r]{1 \\ 1}
            =\colvec{(x+y)/2 \\ (x+y)/2}
          \end{equation*} 
          yang merupakan efek matriks berikut.
          \begin{equation*}
            \begin{mat}[r]
              1/2  &1/2  \\
              1/2  &1/2
            \end{mat}
          \end{equation*}
        \partsitem Memutar seluruh bidang sebesar $\pi/4$~radian searah jarum
          jam membawa garis $y=x$ ke sumbu~$x$.
          Memproyeksikan lalu memutarnya kembali kini memberikan efek yang
          diinginkan.
      \end{exparts}
    \end{answer}
  \item 
    Untuk \( \vec{a},\vec{b}\in\Re^n \), misalkan \( \vec{v}_1 \) adalah
    proyeksi \( \vec{a} \) ke garis yang direntang oleh \( \vec{b} \),
    \( \vec{v}_2 \) adalah proyeksi \( \vec{v}_1 \) ke garis yang direntang
    oleh \( \vec{a} \), \( \vec{v}_3 \) adalah proyeksi \( \vec{v}_2 \) ke
    garis yang direntang oleh \( \vec{b} \), dan seterusnya, bolak-balik antara
    rentang $\vec{a}$ dan $\vec{b}$.
    Artinya, $\vec{v}_{i+1}$ adalah proyeksi $\vec{v}_i$ ke rentang
    $\vec{a}$ jika $i+1$ genap, dan ke rentang $\vec{b}$ jika $i+1$ ganjil.
    Haruskah barisan vektor tersebut akhirnya menetap\Dash yaitu, haruskah ada
    $i$ yang cukup besar sehingga \( \vec{v}_{i+2} \) sama dengan
    \( \vec{v}_{i} \) dan \( \vec{v}_{i+3} \) sama dengan
    \( \vec{v}_{i+1} \)?
    Jika ya, berapakah nilai~$i$ paling awal seperti itu?
    \begin{answer}
      Barisan tersebut tidak harus menetap.
      Dengan
      \begin{equation*}
        \vec{a}=\colvec[r]{1 \\ 0}
        \qquad
        \vec{b}=\colvec[r]{1 \\ 1}
      \end{equation*}
      proyeksi-proyeksinya adalah berikut.
      \begin{equation*}
        \vec{v}_1=\colvec[r]{1/2 \\ 1/2},
        \quad
        \vec{v}_2=\colvec[r]{1/2 \\ 0},
        \quad
        \vec{v}_3=\colvec[r]{1/4 \\ 1/4},
        \quad
        \ldots
      \end{equation*}
      Barisan ini tidak berulang.
      \begin{center}  \small
        \includegraphics{map/mp/ch3.82}
      \end{center}
      \end{answer}
\end{exercises}
















\subsectionoptional{Ortogonalisasi Gram--Schmidt}
\index{proses Gram--Schmidt|(}
% \noindent\textit{This subsection is optional.
%      It is a prerequisite only for the final two sections
%      of Chapter Five.
%      Also, this subsection requires material from the previous subsection, 
%      which itself was optional.}

Subbagian sebelumnya menunjukkan bahwa
memproyeksikan $\vec{v}$ ke garis yang direntang oleh \( \vec{s} \)
menguraikan vektor tersebut menjadi dua bagian
\begin{center}  \small
  \vcenteredhbox{\includegraphics{map/mp/ch3.35}}
  % \begin{tabular}{@{}c@{}}\includegraphics{map/mp/ch3.35}\end{tabular}
   \qquad
   $\displaystyle \vec{v}=\proj{\vec{v}}{\spanof{\vec{s}\,}}
             \,+\,\left(\vec{v}-\proj{\vec{v}}{\spanof{\vec{s}\,}}\right)$
\end{center}
yang ortogonal sehingga ``tidak saling berinteraksi.''
Sekarang kita kembangkan gagasan tersebut.

\begin{definition}  \label{df:MutuallyOrthogonal}
%<*df:MutuallyOrthogonal>
Vektor-vektor \( \vec{v}_1,\dots,\vec{v}_k\in\Re^n \) disebut
\definend{saling ortogonal\/}\index{ortogonal!saling}
jika sebarang dua di antaranya ortogonal:~jika \( i\neq j \), maka
hasil kali titik \( \vec{v}_i\dotprod\vec{v}_j \) adalah nol.
%</df:MutuallyOrthogonal>
\end{definition}

\begin{theorem} \label{th:OrthoIsInd}
%<*th:OrthoIsInd>
Jika vektor-vektor dalam himpunan
\( \set{\vec{v}_1,\dots,\vec{v}_k}\subset\Re^n \) saling ortogonal dan tak
nol, maka himpunan tersebut bebas linear.
%</th:OrthoIsInd>
\end{theorem}

\begin{proof}
%<*pf:OrthoIsInd>
Tinjau
\( \zero=c_1\vec{v}_1+c_2\vec{v}_2+\dots+c_k\vec{v}_k \).
Untuk $i\in \set{1,..\,,k}$, ambil hasil kali titik \( \vec{v}_i \)
dengan kedua ruas persamaan,
$\vec{v}_i\dotprod (c_1\vec{v}_1+c_2\vec{v}_2+\dots+c_k\vec{v}_k)
   =\vec{v}_i\dotprod\zero$,
yang memberikan
$c_i\cdot(\vec{v}_i\dotprod\vec{v}_i)=0$,
Persamaan ini menunjukkan bahwa \( c_i= 0 \) karena
\( \vec{v}_i\neq\zero \).
%</pf:OrthoIsInd>
\end{proof}

\begin{corollary}  \label{cor:OrthAndBigEnoughIsBasis}
%<*co:OrthAndBigEnoughIsBasis>
Dalam ruang vektor berdimensi~$k$, jika vektor-vektor dalam suatu himpunan
berukuran~\( k \) saling ortogonal dan tak nol, maka himpunan tersebut
merupakan basis bagi ruang itu.
%</co:OrthAndBigEnoughIsBasis>
\end{corollary}

\begin{proof}
%<*pf:OrthAndBigEnoughIsBasis>
Setiap subhimpunan bebas linear berukuran~\( k \) dari suatu ruang
berdimensi~$k$ merupakan basis.
%</pf:OrthAndBigEnoughIsBasis>
\end{proof}

Tentu saja, kebalikan \nearbycorollary{cor:OrthAndBigEnoughIsBasis} tidak
berlaku\Dash tidak setiap basis dari setiap subruang $\Re^n$ memiliki
vektor-vektor yang saling ortogonal.
Namun, kita dapat memperoleh kebalikan parsial: bagi setiap subruang $\Re^n$
terdapat setidaknya satu basis yang terdiri atas vektor-vektor yang saling
ortogonal.

\begin{example}
Anggota $\vec{\beta}_1$ dan $\vec{\beta}_2$ dari basis bagi $\Re^2$ berikut
tidak ortogonal.
\begin{center}  \small
  $\displaystyle 
      B=\sequence{
            \colvec[r]{4 \\ 2},
            \colvec[r]{1 \\ 3} }$
  \qquad %\hspace*{3em}
  \begin{tabular}{@{}c@{}}\includegraphics{map/mp/ch3.36}\end{tabular}
\end{center}
Dari $B$ kita akan menurunkan basis baru
$\sequence{\vec{\kappa}_1,\vec{\kappa}_2}$ bagi ruang tersebut, yang terdiri
atas vektor-vektor yang saling ortogonal.
Anggota pertama basis baru itu cukup diambil sebagai $\vec{\beta}_1$.
\begin{equation*}
  \vec{\kappa}_1=\colvec[r]{4 \\ 2}
\end{equation*}
Untuk anggota kedua basis baru, kurangkan dari $\vec{\beta}_2$ bagian yang
searah dengan $\vec{\kappa}_1$.
Yang tersisa adalah bagian $\vec{\beta}_2$ yang ortogonal terhadap
$\vec{\kappa}_1$.
% (it is orthogonal by the definition of the projection into the span of 
% $\vec{\kappa}_1$).
\begin{center}  \small
   $\displaystyle \vec{\kappa}_2=\colvec[r]{1 \\ 3}
    -\proj{\colvec[r]{1 \\ 3}}{\scriptstyle\spanof{\vec{\kappa}_1}}
    =\colvec[r]{1 \\ 3}-\colvec[r]{2 \\ 1}=\colvec[r]{-1 \\ 2}$
   \qquad
  \begin{tabular}{@{}c@{}}\includegraphics{map/mp/ch3.37}\end{tabular}
\end{center} 
Berdasarkan korolari tersebut,
$\sequence{\vec{\kappa}_1,\vec{\kappa}_2}$ merupakan basis bagi $\Re^2$.
\end{example}

\begin{definition}  \label{df:OrthogonalBasis}
%<*df:OrthogonalBasis>
\definend{Basis ortogonal\/}%
\index{ortogonal!basis}\index{basis!ortogonal}
bagi suatu ruang vektor adalah basis yang terdiri atas vektor-vektor yang
saling ortogonal.
%</df:OrthogonalBasis>
\end{definition}

% The next result gives a way to produce an orthogonal basis 
% from any given starting basis.
% We first see an example.

\begin{example} \label{ex:OrthoBasisForReThree}
Untuk menghasilkan basis ortogonal dari basis bagi $\Re^3$ berikut,
\begin{equation*}
  B=\sequence{
           \colvec[r]{1 \\ 1 \\ 1},
           \colvec[r]{0 \\ 2 \\ 0},
           \colvec[r]{1 \\ 0 \\ 3}
            }
\end{equation*}
mulailah dengan mengambil vektor pertama tanpa perubahan.
\begin{equation*}
   \vec{\kappa}_1=\colvec[r]{1 \\ 1 \\ 1}
\end{equation*}
Dapatkan \( \vec{\kappa}_2 \) dengan mengurangkan dari $\vec{\beta}_2$
bagiannya yang searah dengan \( \vec{\kappa}_1 \).
\begin{equation*}
   \vec{\kappa}_2=\colvec[r]{0 \\ 2 \\ 0}
                   -\proj{\colvec[r]{0 \\ 2 \\ 0}}{\spanof{\vec{\kappa}_1}}
                 =\colvec[r]{0 \\ 2 \\ 0}-\colvec[r]{2/3 \\ 2/3 \\ 2/3}
                 =\colvec[r]{-2/3 \\ 4/3 \\ -2/3}
\end{equation*}
Dapatkan $\vec{\kappa}_3$ dengan mengurangkan dari $\vec{\beta}_3$ bagian yang
searah dengan \( \vec{\kappa}_1 \), serta bagian yang searah dengan
\( \vec{\kappa}_2 \).
\begin{equation*}
   \vec{\kappa}_3=\colvec[r]{1 \\ 0 \\ 3}
                   -\proj{\colvec[r]{1 \\ 0 \\ 3}}{\spanof{\vec{\kappa}_1}}
                   -\proj{\colvec[r]{1 \\ 0 \\ 3}}{\spanof{\vec{\kappa}_2}}
                 =\colvec[r]{-1 \\ 0 \\ 1}
\end{equation*}
Seperti di atas, korolari tersebut memberikan bahwa hasilnya merupakan basis
bagi $\Re^3$.
\begin{equation*}
  \sequence{
           \colvec[r]{1 \\ 1 \\ 1},
           \colvec[r]{-2/3 \\ 4/3 \\ -2/3},
           \colvec[r]{-1 \\ 0 \\ 1}
           }
\end{equation*}
\end{example}


\begin{theorem}[Ortogonalisasi Gram--Schmidt]
\index{ortogonalisasi}\index{basis!ortogonalisasi}
\label{th:GramSchmidt}
%<*th:GramSchmidt>
Jika \( \sequence{\vec{\beta}_1,\ldots,\vec{\beta}_k} \) merupakan basis bagi
suatu subruang \( \Re^n \), maka vektor-vektor
\begin{align*}
  \vec{\kappa}_1
  &=\vec{\beta}_1    \\
  \vec{\kappa}_2
  &=\vec{\beta}_2-\proj{\vec{\beta}_2}{\spanof{\vec{\kappa}_1}}    \\
  \vec{\kappa}_3
  &=\vec{\beta}_3
  -\proj{\vec{\beta}_3}{\spanof{\vec{\kappa}_1}}
  -\proj{\vec{\beta}_3}{\spanof{\vec{\kappa}_2}}    \\
  &\vdotswithin{=}           \\
  \vec{\kappa}_k
  &=\vec{\beta}_k
  -\proj{\vec{\beta}_k}{\spanof{\vec{\kappa}_1}}-\cdots
  -\proj{\vec{\beta}_k}{\spanof{\vec{\kappa}_{k-1}}}
\end{align*}
membentuk basis ortogonal bagi subruang yang sama.
%</th:GramSchmidt>
\end{theorem}

\begin{remark}
Hasil ini terbatas pada $\Re^n$ hanya karena kita belum memberikan definisi
keortogonalan bagi ruang-ruang lain.
\end{remark}

\begin{proof}
%<*pf:GramSchmidt0>
Kita akan menggunakan induksi untuk memeriksa bahwa setiap
\( \vec{\kappa}_i \) tak nol, berada dalam rentang
$\sequence{\vec{\beta}_1,\ldots\,\vec{\beta}_i}$, dan ortogonal terhadap semua
vektor sebelumnya,
\( \vec{\kappa}_1\dotprod\vec{\kappa}_i= \cdots
   =\vec{\kappa}_{i-1}\dotprod\vec{\kappa}_{i}=0 \).
Kemudian \nearbycorollary{cor:OrthAndBigEnoughIsBasis} memberikan bahwa
\( \sequence{\vec{\kappa}_1,\ldots\,\vec{\kappa}_k} \) 
merupakan basis bagi ruang yang sama dengan basis awal.
%</pf:GramSchmidt0>

%<*pf:GramSchmidt1>
Kita hanya akan membahas kasus hingga \( i=3 \) untuk memberikan gambaran
tentang argumennya.
Argumen lengkapnya menjadi \nearbyexercise{exer:GramSchmidt}.
%</pf:GramSchmidt1>

%<*pf:GramSchmidt2>
Kasus \( i=1 \) trivial; mengambil \( \vec{\kappa}_1 \) sebagai
\( \vec{\beta}_1 \) menjadikannya vektor tak nol karena $\vec{\beta}_1$
merupakan anggota suatu
basis. Jelas bahwa vektor ini berada dalam rentang
$\sequence{\vec{\beta}_1}$, dan syarat `ortogonal terhadap semua vektor
sebelumnya' terpenuhi secara hampa.
%</pf:GramSchmidt2>

%<*pf:GramSchmidt3>
Dalam kasus \( i=2 \), penguraian
\begin{equation*}
  \vec{\kappa}_2=\vec{\beta}_2
      -\proj{\vec{\beta}_2}{\spanof{\vec{\kappa}_1}}=
  \vec{\beta}_2
  -\frac{\vec{\beta}_2\dotprod\vec{\kappa}_1}{
         \vec{\kappa}_1\dotprod\vec{\kappa}_1}
   \cdot\vec{\kappa}_1
  =
  \vec{\beta}_2
  -\frac{\vec{\beta}_2\dotprod\vec{\kappa}_1}{
         \vec{\kappa}_1\dotprod\vec{\kappa}_1}
   \cdot\vec{\beta}_1
\end{equation*}
menunjukkan bahwa $\vec{\kappa}_2\neq \zero$; jika tidak, persamaan ini akan
menjadi ketergantungan linear tak trivial di antara vektor-vektor
\( \vec{\beta} \) (tak trivial karena koefisien $\vec{\beta}_2$ adalah $1$).
Penguraian tersebut juga menunjukkan bahwa $\vec{\kappa}_2$ berada dalam
rentang $\sequence{\vec{\beta}_1,\vec{\beta}_2}$.
Selain itu, $\vec{\kappa}_2$ ortogonal terhadap satu-satunya vektor sebelumnya,
\begin{equation*}
   \vec{\kappa}_1\dotprod\vec{\kappa}_2=
   \vec{\kappa}_1\dotprod(\vec{\beta}_2
      -\proj{\vec{\beta}_2}{\spanof{\vec{\kappa}_1}})=0
\end{equation*}
karena proyeksi ini ortogonal.
%</pf:GramSchmidt3>

%<*pf:GramSchmidt4>
Kasus \( i=3 \) sama dengan kasus $i=2$, kecuali pada satu perincian.
Seperti dalam kasus $i=2$, uraikan definisinya.
\begin{align*}
  \vec{\kappa}_3
  &=\vec{\beta}_3
  -\frac{\vec{\beta}_3\dotprod\vec{\kappa}_1}{
         \vec{\kappa}_1\dotprod\vec{\kappa}_1}
   \cdot\vec{\kappa}_1
  -\frac{\vec{\beta}_3\dotprod\vec{\kappa}_2}{
         \vec{\kappa}_2\dotprod\vec{\kappa}_2}
   \cdot\vec{\kappa}_2  \\
  &=\vec{\beta}_3
  -\frac{\vec{\beta}_3\dotprod\vec{\kappa}_1}{
         \vec{\kappa}_1\dotprod\vec{\kappa}_1}
   \cdot\vec{\beta}_1
  -\frac{\vec{\beta}_3\dotprod\vec{\kappa}_2}{
         \vec{\kappa}_2\dotprod\vec{\kappa}_2}
  \cdot\bigl(\vec{\beta}_2
  -\frac{\vec{\beta}_2\dotprod\vec{\kappa}_1}{
         \vec{\kappa}_1\dotprod\vec{\kappa}_1}
   \cdot\vec{\beta}_1\bigr)
\end{align*}
Berdasarkan baris pertama, $\vec{\kappa}_3\neq\zero$, karena
$\vec{\beta}_3$ tidak berada dalam rentang
$\spanof{\vec{\beta}_1,\vec{\beta}_2}$ dan oleh karena itu, berdasarkan
hipotesis induksi, tidak berada dalam rentang
$\spanof{\vec{\kappa}_1,\vec{\kappa}_2}$.
Berdasarkan baris kedua, $\vec{\kappa}_3$ berada dalam rentang tiga
vektor~$\vec{\beta}$ pertama.
Akhirnya, perhitungan berikut menunjukkan bahwa $\vec{\kappa}_3$ ortogonal
terhadap $\vec{\kappa}_1$.
%</pf:GramSchmidt4>
%<*pf:GramSchmidt5>
\begin{align*}
   \vec{\kappa}_1\dotprod\vec{\kappa}_3
   &=\vec{\kappa}_1\dotprod\bigl(\;\vec{\beta}_3
      -\proj{\vec{\beta}_3}{\spanof{\vec{\kappa}_1}}
      -\proj{\vec{\beta}_3}{\spanof{\vec{\kappa}_2}}\;\bigr) \\
   &=\vec{\kappa}_1\dotprod\bigl(\vec{\beta}_3
          -\proj{\vec{\beta}_3}{\spanof{\vec{\kappa}_1}}\bigr)
    -\vec{\kappa}_1\dotprod\proj{\vec{\beta}_3}{\spanof{\vec{\kappa}_2}} \\
   &=0
\end{align*}
(Inilah perbedaannya dengan kasus $i=2$:
seperti pada $i=2$, suku pertama adalah~$0$ karena proyeksi ini ortogonal,
tetapi di sini suku kedua pada baris kedua adalah~$0$ karena
\( \vec{\kappa}_1 \) ortogonal terhadap \( \vec{\kappa}_2 \), sehingga juga
ortogonal terhadap sebarang vektor pada garis yang direntang oleh
\( \vec{\kappa}_2 \).)
Pemeriksaan serupa menunjukkan bahwa $\vec{\kappa}_3$ juga ortogonal terhadap
$\vec{\kappa}_2$.
%</pf:GramSchmidt5>
\end{proof}

Selain membuat vektor-vektor dalam basis saling ortogonal, kita juga dapat
\definend{menormalisasi}\index{normalisasi, vektor} setiap vektor dengan
membaginya dengan panjangnya, sehingga akhirnya diperoleh
\definend{basis ortonormal}.\index{basis!ortonormal}%
\index{basis ortonormal}.

\begin{example}
Dari basis ortogonal dalam \nearbyexample{ex:OrthoBasisForReThree}, normalisasi
menghasilkan basis ortonormal berikut.
\begin{equation*}
  \sequence{
           \colvec[r]{1/\sqrt{3} \\ 1/\sqrt{3} \\ 1/\sqrt{3}},
           \colvec[r]{-1/\sqrt{6} \\ 2/\sqrt{6} \\ -1/\sqrt{6}},
           \colvec[r]{-1/\sqrt{2} \\ 0 \\ 1/\sqrt{2}}
           }
\end{equation*}
\end{example}

\noindent Selain menarik secara intuitif dan serupa dengan basis standar
$\stdbasis_n$ bagi $\Re^n$, basis ortonormal juga menyederhanakan beberapa
perhitungan.
\nearbyexercise{exer:OrthoRepEasy} merupakan salah satu contohnya.




\begin{exercises}
  \item Normalisasikan panjang vektor-vektor berikut.
    \begin{exparts*}
      \partsitem $\colvec{1 \\ 2}$
      \partsitem $\colvec{-1 \\ 3 \\ 0}$
      \partsitem $\colvec{1 \\ -1}$
    \end{exparts*}
    \begin{answer}
      \begin{exparts}
        \partsitem $\colvec{1/\sqrt{5} \\ 2/\sqrt{5}}
             =\frac{1}{\sqrt{5}}\cdot\colvec{1 \\ 2}$
        \partsitem $\frac{1}{\sqrt{10}}\cdot\colvec{-1 \\ 3 \\ 0}$
        \partsitem $\frac{1}{\sqrt{2}}\cdot\colvec{1 \\ -1}$
      \end{exparts}
    \end{answer}
  \recommended \item Terapkan Gram--Schmidt pada basis bagi~$\Re^2$ berikut.
    \begin{equation*}
      \sequence{\colvec{1 \\ 1}, \colvec{-1 \\ 2}}
    \end{equation*}
    Periksa bahwa vektor-vektor yang dihasilkan saling ortogonal.
    \begin{answer}
      Namai basis yang diberikan
      $B=\sequence{\vec{\beta}_1,\vec{\beta}_2}$.
      Pertama, $\vec{\kappa}_1=\vec{\beta}_1$.
      Untuk vektor lainnya,
      $\vec{\kappa}_2=\vec{\beta}_2-\proj{\vec{\beta}_2}{\spanof{\vec{\kappa}_1}}$.
      \begin{equation*}
        \vec{\kappa}_2=
        \colvec{-1 \\ 2}-\frac{\colvec{-1 \\ 2}\dotprod\colvec{1 \\ 1}}{\colvec{1 \\ 1}\dotprod\colvec{1 \\ 1}}\cdot\colvec{1 \\ 1}
        =\colvec{-3/2 \\ 3/2}
      \end{equation*}
      Untuk memeriksa bahwa keduanya ortogonal, cukup perhatikan bahwa hasil
      kali titik keduanya nol.
    \end{answer}
  \recommended \item 
    Terapkan proses Gram--Schmidt pada basis bagi $\Re^3$ berikut.
    \begin{equation*}
      \sequence{
                         \colvec{1 \\ 2 \\ 3},
                         \colvec{2 \\ 1 \\ -3},
                         \colvec{3 \\ 3 \\ 3}
                         }  
    \end{equation*}
    \begin{answer}
      Namai basis yang diberikan
      $B=\sequence{\vec{\beta}_1,\vec{\beta}_2,\vec{\beta}_3}$.
      Pertama, $\vec{\kappa}_1=\vec{\beta}_1$.

      Selanjutnya,
      $\vec{\kappa}_2
          =\vec{\beta}_2-\proj{\vec{\beta}_2}{\spanof{\vec{\kappa}_1}}$.
      \begin{equation*}
        \vec{\kappa}_2=
        \colvec{2 \\ 1 \\ -3}-\frac{\colvec{2 \\ 1 \\ -3}\dotprod\colvec{1 \\ 2 \\ 3}}{\colvec{1 \\ 2 \\ 3}\dotprod\colvec{1 \\ 2 \\ 3}}\cdot\colvec{1 \\ 2 \\ 3}
        =\colvec{33/14 \\ 24/14 \\ -27/14}
      \end{equation*}

      Untuk vektor ketiga, rumusnya adalah
      \begin{equation*}
         \vec{\kappa}_3
          =\vec{\beta}_3
           -\proj{\vec{\beta}_3}{\spanof{\vec{\kappa}_1}}
           -\proj{\vec{\beta}_3}{\spanof{\vec{\kappa}_2}}    
      \end{equation*} 
      dan berikut perhitungannya.
      \begin{equation*}
        \colvec{3 \\ 3 \\ 3}
        -\frac{\colvec{3 \\ 3 \\ 3}\dotprod\colvec{1 \\ 2 \\ 3}}{\colvec{1 \\ 2 \\ 3}\dotprod\colvec{1 \\ 2 \\ 3}}\cdot\colvec{1 \\ 2 \\ 3}
        -\frac{\colvec{3 \\ 3 \\ 3}\dotprod\colvec{33/14 \\ 24/14 \\ -27/14}}{\colvec{33/14 \\ 24/14 \\ -27/14}\dotprod\colvec{33/14 \\ 24/14 \\ -27/14}}\cdot\colvec{33/14 \\ 24/14 \\ -27/14}
       =\colvec{9/19 \\ -9/19 \\ 3/19}
      \end{equation*}
      % 1 - 18/14-(0/14)/(1274/196) times \colvec{3 \\ 3 \\ 3}
    \end{answer}
  \recommended \item 
    Terapkan Gram--Schmidt pada setiap basis bagi $\Re^2$ berikut.
    \begin{exparts*}
      \partsitem \( \sequence{
                         \colvec[r]{1 \\ 1},
                         \colvec[r]{2 \\ 1}
                         }  \)
      \partsitem \( \sequence{
                         \colvec[r]{0 \\ 1},
                         \colvec[r]{-1 \\ 3}
                         }  \)
      \partsitem \( \sequence{
                         \colvec[r]{0 \\ 1},
                         \colvec[r]{-1 \\ 0}
                         }  \)
    \end{exparts*}
    Kemudian ubah basis-basis ortogonal tersebut menjadi basis ortonormal.
    \begin{answer}
      \begin{exparts}
       \partsitem 
        \begin{align*}
          \vec{\kappa}_1 &= \colvec[r]{1 \\ 1}           \\
          \vec{\kappa}_2
            &=
            \colvec[r]{2 \\ 1}
            -\proj{\colvec[r]{2 \\ 1}}{\spanof{\vec{\kappa}_1}}  
            =
            \colvec[r]{2 \\ 1}
            -\frac{\colvec[r]{2 \\ 1}\dotprod\colvec[r]{1 \\ 1}}{%
                    \colvec[r]{1 \\ 1}\dotprod\colvec[r]{1 \\ 1}}
            \cdot\colvec[r]{1 \\ 1}                                \\
            &\quad=
            \colvec[r]{2 \\ 1}
            -\frac{3}{2}
            \cdot\colvec[r]{1 \\ 1}                                
            =\colvec[r]{1/2 \\ -1/2}
        \end{align*}
        Berikut basis ortonormal yang bersesuaian.
        \begin{equation*}
          \sequence{
              \colvec[r]{1/\sqrt{2} \\ 1/\sqrt{2}},
              \colvec[r]{\sqrt{2}/2 \\ -\sqrt{2}/2}
                }
        \end{equation*}
       \partsitem 
        \begin{align*}
          \vec{\kappa}_1 &= \colvec[r]{0 \\ 1}           \\
          \vec{\kappa}_2
            &=
            \colvec[r]{-1 \\ 3}
            -\proj{\colvec[r]{-1 \\ 3}}{\spanof{\vec{\kappa}_1}}  
            =
            \colvec[r]{-1 \\ 3}
            -\frac{\colvec[r]{-1 \\ 3}\dotprod\colvec[r]{0 \\ 1}}{%
                    \colvec[r]{0 \\ 1}\dotprod\colvec[r]{0 \\ 1}}
            \cdot\colvec[r]{0 \\ 1}                                \\
            &\quad=
            \colvec[r]{-1 \\ 3}
            -\frac{3}{1}
            \cdot\colvec[r]{0 \\ 1}                              
            =\colvec[r]{-1 \\ 0}
        \end{align*}
        Berikut basis ortonormalnya.
        \begin{equation*}
          \sequence{
                \colvec[r]{0 \\ 1},
                \colvec[r]{-1 \\ 0}
                }
        \end{equation*}
       \partsitem 
        \begin{align*}
          \vec{\kappa}_1 &= \colvec[r]{0 \\ 1}           \\
          \vec{\kappa}_2
            &=
            \colvec[r]{-1 \\ 0}
            -\proj{\colvec[r]{-1 \\ 0}}{\spanof{\vec{\kappa}_1}}  
            =
            \colvec[r]{-1 \\ 0}
            -\frac{\colvec[r]{-1 \\ 0}\dotprod\colvec[r]{0 \\ 1}}{%
                    \colvec[r]{0 \\ 1}\dotprod\colvec[r]{0 \\ 1}}
            \cdot\colvec[r]{0 \\ 1}                               \\ 
            &\quad=
            \colvec[r]{-1 \\ 0}
            -\frac{0}{1}
            \cdot\colvec[r]{0 \\ 1}                                
            =\colvec[r]{-1 \\ 0}
        \end{align*}
      \end{exparts}
      Berikut basis ortonormal yang terkait.
      \begin{equation*}
        \sequence{
              \colvec[r]{0 \\ 1},
              \colvec[r]{-1 \\ 0}
              }
      \end{equation*}
    \end{answer}
  \item 
    Terapkan proses Gram--Schmidt pada setiap basis bagi $\Re^3$ berikut.
    \begin{exparts*}
      \partsitem \( \sequence{
                         \colvec[r]{2 \\ 2 \\ 2},
                         \colvec[r]{1 \\ 0 \\ -1},
                         \colvec[r]{0 \\ 3 \\ 1}
                         }  \)
      \partsitem \( \sequence{
                         \colvec[r]{1 \\ -1 \\ 0},
                         \colvec[r]{0 \\ 1 \\ 0},
                         \colvec[r]{2 \\ 3 \\ 1}
                         }  \)
    \end{exparts*}
    Kemudian ubah basis-basis ortogonal tersebut menjadi basis ortonormal.
    \begin{answer} 
      \begin{exparts}
       \partsitem Vektor basis pertama tidak berubah.
        \begin{equation*}
          \vec{\kappa}_1 = \colvec[r]{2 \\ 2 \\ 2}          
        \end{equation*}
        Vektor kedua diperoleh dari perhitungan berikut.
        \begin{align*}
          \vec{\kappa}_2
            &=
            \colvec[r]{1 \\ 0 \\ -1}
            -\proj{\colvec[r]{1 \\ 0 \\ -1}}{\spanof{\vec{\kappa}_1}}  
            =
            \colvec[r]{1 \\ 0 \\ -1}
            -\frac{\colvec[r]{1 \\ 0 \\ -1}\dotprod\colvec[r]{2 \\ 2 \\ 2}}{%
                    \colvec[r]{2 \\ 2 \\ 2}\dotprod\colvec[r]{2 \\ 2 \\ 2}}
            \cdot\colvec[r]{2 \\ 2 \\ 2}                                   \\
            &=
            \colvec[r]{1 \\ 0 \\ -1}
            -\frac{0}{12}
            \cdot\colvec[r]{2 \\ 2 \\ 2}                                
            =\colvec[r]{1 \\ 0 \\ -1}                              
        \end{align*}
        Untuk vektor ketiga, aritmetikanya lebih rumit tetapi perhitungannya
        langsung.
        \begin{align*}
          \vec{\kappa}_3
            &=
            \colvec[r]{0 \\ 3 \\ 1}
            -\proj{\colvec[r]{0 \\ 3 \\ 1}}{\spanof{\vec{\kappa}_1}}  
            -\proj{\colvec[r]{0 \\ 3 \\ 1}}{\spanof{\vec{\kappa}_2}}   \\ 
            &=
            \colvec[r]{0 \\ 3 \\ 1}
            -\frac{\colvec[r]{0 \\ 3 \\ 1}\dotprod\colvec[r]{2 \\ 2 \\ 2}}{%
                    \colvec[r]{2 \\ 2 \\ 2}\dotprod\colvec[r]{2 \\ 2 \\ 2}}
            \cdot\colvec[r]{2 \\ 2 \\ 2}                                
            -\frac{\colvec[r]{0 \\ 3 \\ 1}\dotprod\colvec[r]{1 \\ 0 \\ -1}}{%
                    \colvec[r]{1 \\ 0 \\ -1}\dotprod\colvec[r]{1 \\ 0 \\ -1}}
            \cdot\colvec[r]{1 \\ 0 \\ -1}                                   \\
            &\hbox{}\quad=
            \colvec[r]{0 \\ 3 \\ 1}
            -\frac{8}{12}
            \cdot\colvec[r]{2 \\ 2 \\ 2}                                
            -\frac{-1}{2}
            \cdot\colvec[r]{1 \\ 0 \\ -1}                              
            =\colvec[r]{-5/6 \\ 5/3 \\ -5/6}
        \end{align*}
        Berikut basis ortonormalnya.
        \begin{equation*}
          \sequence{
              \colvec[r]{1/\sqrt{3} \\ 1/\sqrt{3} \\ 1/\sqrt{3}},
              \colvec[r]{1/\sqrt{2} \\ 0 \\ -1/\sqrt{2}},
              \colvec[r]{-1/\sqrt{6} \\ 2/\sqrt{6} \\ -1/\sqrt{6}}
                }
        \end{equation*}
       \partsitem Vektor basis pertama adalah vektor yang diberikan.
        \begin{equation*}
          \vec{\kappa}_1 = \colvec[r]{1 \\ -1 \\ 0}           
        \end{equation*}
        Berikut vektor keduanya.
        \begin{align*}
          \vec{\kappa}_2
            &=
            \colvec[r]{0 \\ 1 \\ 0}
            -\proj{\colvec[r]{0 \\ 1 \\ 0}}{\spanof{\vec{\kappa}_1}}  
            =
            \colvec[r]{0 \\ 1 \\ 0}
            -\frac{\colvec[r]{0 \\ 1 \\ 0}\dotprod\colvec[r]{1 \\ -1 \\ 0}}{%
                    \colvec[r]{1 \\ -1 \\ 0}\dotprod\colvec[r]{1 \\ -1 \\ 0}}
            \cdot\colvec[r]{1 \\ -1 \\ 0}                                   \\ 
            &=
            \colvec[r]{0 \\ 1 \\ 0}
            -\frac{-1}{2}
            \cdot\colvec[r]{1 \\ -1 \\ 0}                                
            =\colvec[r]{1/2 \\ 1/2 \\ 0}                                  
         \end{align*}
         Berikut vektor ketiganya.
        \begin{align*}
          \vec{\kappa}_3
            &=
            \colvec[r]{2 \\ 3 \\ 1}
            -\proj{\colvec[r]{2 \\ 3 \\ 1}}{\spanof{\vec{\kappa}_1}}  
            -\proj{\colvec[r]{2 \\ 3 \\ 1}}{\spanof{\vec{\kappa}_2}}       \\ 
            &=
            \colvec[r]{2 \\ 3 \\ 1}
            -\frac{\colvec[r]{2 \\ 3 \\ 1}\dotprod\colvec[r]{1 \\ -1 \\ 0}}{%
                    \colvec[r]{1 \\ -1 \\ 0}\dotprod\colvec[r]{1 \\ -1 \\ 0}}
            \cdot\colvec[r]{1 \\ -1 \\ 0}                                
            -\frac{\colvec[r]{2 \\ 3 \\ 1}\dotprod\colvec[r]{1/2 \\ 1/2 \\ 0}}{%
                    \colvec[r]{1/2 \\ 1/2 \\ 0}\dotprod\colvec[r]{1/2 \\ 1/2 \\ 0}}
            \cdot\colvec[r]{1/2 \\ 1/2 \\ 0}                   \\
            &=
            \colvec[r]{2 \\ 3 \\ 1}
            -\frac{-1}{2}
            \cdot\colvec[r]{1 \\ -1 \\ 0}                                
            -\frac{5/2}{1/2}
            \cdot\colvec[r]{1/2 \\ 1/2 \\ 0} 
            =\colvec[r]{0 \\ 0 \\ 1}
        \end{align*}
        Berikut basis ortonormal yang terkait.
        \begin{equation*}
          \sequence{
              \colvec[r]{1/\sqrt{2} \\ -1/\sqrt{2} \\ 0},
              \colvec[r]{1/\sqrt{2} \\ 1/\sqrt{2} \\ 0},
              \colvec[r]{0 \\ 0 \\ 1}
                }
        \end{equation*}
      \end{exparts}
    \end{answer}
   \recommended \item 
       Temukan basis ortonormal bagi subruang $\Re^3$ berikut:~bidang
       $x-y+z=0$.
       \begin{answer}
         Kita dapat memparameterkan ruang yang diberikan dengan cara berikut.
         \begin{equation*}
           \set{\colvec{x \\ y \\ z} \suchthat x=y-z}
           =\set{\colvec[r]{1 \\ 1 \\ 0}\cdot y+\colvec[r]{-1 \\ 0 \\ 1}\cdot z 
                     \suchthat y,z\in\Re}
         \end{equation*}
         Jadi, kita ambil basis
         \begin{equation*}
           \sequence{\colvec[r]{1 \\ 1 \\ 0},
                     \colvec[r]{-1 \\ 0 \\ 1}}
         \end{equation*}
         lalu terapkan proses Gram--Schmidt untuk memperoleh vektor basis
         pertama berikut
         \begin{equation*}
           \vec{\kappa}_1 = \colvec[r]{1 \\ 1 \\ 0}             
         \end{equation*}
         dan vektor kedua berikut.
         \begin{align*}
           \vec{\kappa}_2
           &=
           \colvec[r]{-1 \\ 0 \\ 1}-\proj{\colvec[r]{-1 \\ 0 \\ 1}}{%
                                        \spanof{\vec{\kappa}_1}} 
           =
           \colvec[r]{-1 \\ 0 \\ 1}
            -\frac{\colvec[r]{-1 \\ 0 \\ 1}\dotprod\colvec[r]{1 \\ 1 \\ 0}}{%
                    \colvec[r]{1 \\ 1 \\ 0}\dotprod\colvec[r]{1 \\ 1 \\ 0}}
              \cdot\colvec[r]{1 \\ 1 \\ 0}                                \\
           &=
           \colvec[r]{-1 \\ 0 \\ 1}
            -\frac{-1}{2}\cdot\colvec[r]{1 \\ 1 \\ 0}                  
           =\colvec[r]{-1/2 \\ 1/2 \\ 1}
         \end{align*}
         lalu lakukan normalisasi.
         \begin{equation*}
           \sequence{
                    \colvec[r]{1/\sqrt{2} \\ 1/\sqrt{2} \\ 0},
                    \colvec[r]{-1/\sqrt{6} \\ 1/\sqrt{6} \\ 2/\sqrt{6}}
                    }
         \end{equation*}
       \end{answer}
   \item 
     Temukan basis ortonormal bagi subruang $\Re^4$ berikut.
     \begin{equation*}
       \set{\colvec{x \\ y \\ z \\ w}\suchthat x-y-z+w=0\text{\ dan\ }x+z=0}
     \end{equation*}
     \begin{answer}
       Mereduksi sistem linear
       \begin{equation*}
         \begin{linsys}{4}
           x  &-  &y  &-  &z  &+  &w  &=  &0  \\
           x  &   &   &+  &z  &   &   &=  &0
         \end{linsys}
         \grstep{-\rho_1+\rho_2}
         \begin{linsys}{4}
           x  &-  &y  &-  &z  &+  &w  &=  &0  \\
              &   &y  &+  &2z &-  &w  &=  &0
         \end{linsys}
       \end{equation*}
       dan memparameterkannya memberikan deskripsi subruang berikut.
       \begin{equation*}
         \set{\colvec[r]{-1 \\ -2 \\ 1 \\ 0}\cdot z
              +\colvec[r]{0 \\ 1 \\ 0 \\ 1}\cdot w
              \suchthat z,w\in\Re}
       \end{equation*}
       Jadi, kita ambil basis
       \begin{equation*}
         \sequence{
                  \colvec[r]{-1 \\ -2 \\ 1 \\ 0},
                  \colvec[r]{0 \\ 1 \\ 0 \\ 1}
                  }
       \end{equation*}
       lalu menjalankan proses Gram--Schmidt dengan vektor basis pertama
       \begin{equation*}
         \vec{\kappa}_1 = \colvec[r]{-1 \\ -2 \\ 1 \\ 0}   
       \end{equation*}
       dan vektor basis kedua
       \begin{align*}
         \vec{\kappa}_2
          &=
          \colvec[r]{0 \\ 1 \\ 0 \\ 1}
          -\proj{\colvec[r]{0 \\ 1 \\ 0 \\ 1}}{\spanof{\vec{\kappa}_1}}   
          =
          \colvec[r]{0 \\ 1 \\ 0 \\ 1}
         -\frac{\colvec[r]{0 \\ 1 \\ 0 \\ 1}\dotprod\colvec[r]{-1 \\ -2 \\ 1 \\ 0}}{%
             \colvec[r]{-1 \\ -2 \\ 1 \\ 0}\dotprod\colvec[r]{-1 \\ -2 \\ 1 \\ 0}} 
          \cdot\colvec[r]{-1 \\ -2 \\ 1 \\ 0}                            \\
          &=
          \colvec[r]{0 \\ 1 \\ 0 \\ 1}
         -\frac{-2}{6} 
          \cdot\colvec[r]{-1 \\ -2 \\ 1 \\ 0}  
          =\colvec[r]{-1/3 \\ 1/3 \\ 1/3 \\ 1}
       \end{align*}
       dan mengakhirinya dengan normalisasi.
       \begin{equation*}
         \sequence{
               \colvec[r]{-1/\sqrt{6} \\ -2/\sqrt{6} \\ 1/\sqrt{6} \\ 0},
               \colvec[r]{-\sqrt{3}/6 \\ \sqrt{3}/6 \\ \sqrt{3}/6 \\ \sqrt{3}/2}
                  }
       \end{equation*}
     \end{answer}
  \item 
     Tunjukkan bahwa sebarang subhimpunan bebas linear dari \( \Re^n \) dapat
     diortogonalisasi tanpa mengubah rentangnya.
     \begin{answer} 
       Suatu subhimpunan bebas linear dari $\Re^n$ merupakan basis bagi
       rentangnya sendiri.
       Terapkan \nearbytheorem{th:GramSchmidt}.

       \textit{Catatan}.
       Berikut alasan frasa `bebas linear' dicantumkan dalam soal.
       Menghilangkan frasa tersebut mengharuskan kita memperhatikan dua hal.
       Hal pertama ialah bahwa ketika kita menjalankan proses Gram--Schmidt
       pada suatu himpunan yang bergantung linear, kita memperoleh beberapa
       vektor nol.
       Sebagai contoh, dengan
       \begin{equation*}
         S=\set{\colvec[r]{1 \\ 2},\colvec[r]{3 \\ 6}}
       \end{equation*}
       kita akan memperoleh
       \begin{equation*}
         \vec{\kappa}_1   =   \colvec[r]{1 \\ 2}
         \qquad 
         \vec{\kappa}_2
           =
           \colvec[r]{3 \\ 6}-\proj{\colvec[r]{3 \\ 6}}{\spanof{\vec{\kappa}_1}} 
           =
           \colvec[r]{0 \\ 0}
       \end{equation*}
       Hal pertama ini tidak terlalu buruk karena vektor nol menurut definisi
       ortogonal terhadap setiap vektor lain. Jadi, kita dapat menerima situasi
       ini sebagai penghasil himpunan ortogonal (meskipun tentu tidak dapat
       dinormalisasi), atau cukup mengubah prosedur Gram--Schmidt agar membuang
       setiap vektor nol.
       Hal kedua yang perlu diperhatikan jika frasa `bebas linear' dihilangkan
       dari soal adalah bahwa himpunannya mungkin tak hingga.
       Tentu saja, setiap subruang dari $\Re^n$ yang berdimensi hingga juga
       berdimensi hingga, sehingga hanya berhingga banyak anggotanya yang dapat
       bebas linear. Meskipun demikian, suatu ``proses'' yang memeriksa satu
       demi satu vektor dalam himpunan tak hingga setidaknya memerlukan
       penjelasan tambahan dalam soal ini.
       Subhimpunan bebas linear dari $\Re^n$ otomatis berhingga\Dash bahkan
       berukuran paling besar~$n$\Dash sehingga frasa `bebas linear'
       meniadakan persoalan-persoalan tersebut.
     \end{answer}
  \item Apa yang terjadi jika kita mencoba menerapkan proses Gram--Schmidt pada
    suatu himpunan berhingga yang bukan basis?
    \begin{answer}
      Jika himpunan itu tidak bebas linear, kita memperoleh vektor nol.
      Jika tidak (jika himpunan kita bebas linear tetapi tidak merentang
      ruang), kita menjalankan Gram--Schmidt pada himpunan yang merupakan basis
      bagi suatu subruang, sehingga kita memperoleh basis ortogonal bagi
      subruang tersebut.
    \end{answer}
   \recommended \item
     Apa yang terjadi jika kita menerapkan proses Gram--Schmidt pada basis yang
     sudah ortogonal?
     \begin{answer}
       Proses tersebut tidak mengubah basis.
     \end{answer}
  \item 
     Misalkan $\sequence{\vec{\kappa}_1,\dots,\vec{\kappa}_k}$ merupakan
     himpunan vektor-vektor yang saling ortogonal dalam $\Re^n$.
     \begin{exparts}
       \partsitem Buktikan bahwa bagi sebarang $\vec{v}$ dalam ruang tersebut,
         vektor
         $\vec{v}-(\proj{\vec{v}\,}{\spanof{\vec{\kappa}_1}}
           +\dots+\proj{\vec{v}\,}{\spanof{\vec{\kappa}_k}})$
         ortogonal terhadap setiap $\vec{\kappa}_1$, \ldots,
         $\vec{\kappa}_k$.
       \partsitem Ilustrasikan bagian sebelumnya dalam $\Re^3$ dengan
         menggunakan $\vec{e}_1$ sebagai $\vec{\kappa}_1$, $\vec{e}_2$ sebagai
         $\vec{\kappa}_2$, dan mengambil $\vec{v}$ dengan komponen $1$, $2$,
         dan $3$.
       \partsitem Tunjukkan bahwa
         $\proj{\vec{v}\,}{\spanof{\vec{\kappa}_1}}
         +\dots+\proj{\vec{v}\,}{\spanof{\vec{\kappa}_k}}$ adalah vektor
         dalam rentang himpunan vektor-vektor $\vec{\kappa}$ yang paling dekat
         dengan $\vec{v}$.
         \textit{Petunjuk}. Pada ilustrasi yang dibuat untuk bagian sebelumnya,
         tambahkan vektor $d_1\vec{\kappa}_1+d_2\vec{\kappa}_2$, lalu terapkan
         Teorema Pythagoras pada segitiga yang dihasilkan.
     \end{exparts}
     \begin{answer}
       \begin{exparts}
         \partsitem Argumennya sama seperti kasus $i=3$ dalam bukti
           \nearbytheorem{th:GramSchmidt}.
           Hasil kali titik
           \begin{equation*}
             \vec{\kappa}_i\dotprod
             \left(\vec{v}-\proj{\vec{v}\,}{\spanof{\vec{\kappa}_1}}
               -\dots-\proj{\vec{v}\,}{\spanof{\vec{v}_k}}\right)
           \end{equation*}
           dapat dituliskan sebagai jumlah suku-suku berbentuk
           $-\vec{\kappa}_i\dotprod\proj{\vec{v}\,}{\spanof{\vec{\kappa}_j}}$ 
           dengan $j\neq i$, serta suku
           $\vec{\kappa}_i\dotprod(\vec{v}-
                                 \proj{\vec{v}\,}{\spanof{\vec{\kappa}_i}})$.
           Suku jenis pertama sama dengan nol karena vektor-vektor
           $\vec{\kappa}$ saling ortogonal.
           Suku lainnya nol karena proyeksi ini ortogonal
           (yakni, definisi proyeksi membuatnya nol:
           $\vec{\kappa}_i\dotprod(\vec{v}
               -\proj{\vec{v}\,}{\spanof{\vec{\kappa}_i}})
            =\vec{\kappa}_i\dotprod\vec{v}-
              \vec{\kappa}_i\dotprod
               \left((\vec{v}\dotprod\vec{\kappa}_i)/(%
                      \vec{\kappa}_i\dotprod\vec{\kappa}_i)\right)
                   \cdot\vec{\kappa}_i$
           setelah semua pencoretan dilakukan, hasilnya sama dengan nol).
         \partsitem Vektor $\vec{v}$ berwarna hitam, sedangkan vektor
           vektor $\proj{\vec{v}\,}{\spanof{\vec{\kappa}_1}}
                    +\proj{\vec{v}\,}{\spanof{\vec{\kappa}_2}}
                   =1\cdot\vec{e}_1+2\cdot\vec{e}_2$ berwarna abu-abu.
           \begin{center}  \small
             \includegraphics{map/mp/ch3.83}
              \end{center}
              Vektor
              $\vec{v}-(\proj{\vec{v}\,}{\spanof{\vec{\kappa}_1}}
                    +\proj{\vec{v}\,}{\spanof{\vec{\kappa}_2}})$
              terletak pada garis putus-putus yang menghubungkan vektor hitam
              dengan vektor abu-abu; artinya, vektor itu ortogonal terhadap
              bidang~$xy$.
          \partsitem Dengan mengikuti petunjuk, kita memperoleh diagram berikut.
           \begin{center}  \small
             \includegraphics{map/mp/ch3.84}
            \end{center}
            Segitiga putus-putus bersudut siku-siku pada tempat vektor abu-abu
            $1\cdot\vec{e}_1+2\cdot\vec{e}_2$ bertemu garis putus-putus vertikal
            $\vec{v}-(1\cdot\vec{e}_1+2\cdot\vec{e}_2)$; inilah yang dibuktikan
            oleh bagian pertama soal ini.
            Teorema Pythagoras kemudian memberikan bahwa sisi miring\Dash ruas
            dari $\vec{v}$ ke sebarang vektor lain\Dash setidaknya sepanjang
            garis putus-putus vertikal.
       
            Secara lebih formal, tuliskan
            $\proj{\vec{v}\,}{\spanof{\vec{\kappa}_1}}
                 +\dots+\proj{\vec{v}\,}{\spanof{\vec{\kappa}_k}}$ sebagai
            $c_1\cdot\vec{\kappa}_1+\dots+c_k\cdot\vec{\kappa}_k$, lalu tinjau
            sebarang vektor lain dalam rentang tersebut,
            $d_1\cdot\vec{\kappa}_1+\dots+d_k\cdot\vec{\kappa}_k$. 
            Perhatikan bahwa
            \begin{multline*}
              \vec{v}-(d_1\cdot\vec{\kappa}_1+\dots+d_k\cdot\vec{\kappa}_k)  \\
              \begin{aligned}
              &=
              \bigl(\vec{v}
                    -(c_1\cdot\vec{\kappa}_1+\dots+c_k\cdot\vec{\kappa}_k)
              \bigr)                                                     \\
              &\quad+\bigl((c_1\cdot\vec{\kappa}_1+\dots+c_k\cdot\vec{\kappa}_k)
                     -(d_1\cdot\vec{\kappa}_1+\dots+d_k\cdot\vec{\kappa}_k)
               \bigr)
              \end{aligned}
            \end{multline*}
            dan bahwa
            $
              \bigl(\vec{v}
                 -(c_1\cdot\vec{\kappa}_1+\dots+c_k\cdot\vec{\kappa}_k)
              \bigr)
              \dotprod
              \bigl((c_1\cdot\vec{\kappa}_1+\dots+c_k\cdot\vec{\kappa}_k)
                     -(d_1\cdot\vec{\kappa}_1+\dots+d_k\cdot\vec{\kappa}_k)
              \bigr)                                                      
              =0
            $
            (karena bagian pertama menunjukkan bahwa $\vec{v}
                 -(c_1\cdot\vec{\kappa}_1+\dots+c_k\cdot\vec{\kappa}_k)$
            ortogonal terhadap setiap $\vec{\kappa}$ sehingga juga ortogonal
            terhadap kombinasi linear vektor-vektor $\vec{\kappa}$ ini).
            Sekarang terapkan Teorema Pythagoras.
       \end{exparts}
     \end{answer}
   \item 
     Temukan vektor tak nol dalam \( \Re^3 \) yang ortogonal terhadap kedua
     vektor berikut.
     \begin{equation*}
       \colvec[r]{1 \\ 5 \\ -1}
       \qquad
       \colvec[r]{2 \\ 2 \\ 0}
     \end{equation*}
     \begin{answer} 
       Salah satu caranya ialah mencari vektor ketiga sehingga ketiganya
       bersama-sama membentuk basis bagi $\Re^3$, misalnya,
       \begin{equation*}
         \vec{\beta}_3=\colvec[r]{1 \\ 0 \\ 0}  
       \end{equation*} 
       (vektor kedua tidak bergantung pada vektor ketiga karena memiliki
       komponen kedua yang tak nol, dan vektor pertama tidak bergantung pada
       vektor kedua dan ketiga karena komponen ketiganya tak nol), lalu terapkan
       proses Gram--Schmidt.
       Anggota pertama basis baru adalah berikut.
       \begin{equation*}
         \vec{\kappa}_1 = \colvec[r]{1 \\ 5 \\ -1}  
       \end{equation*}
       Berikut anggota keduanya.
       \begin{align*}
         \vec{\kappa}_2
           &=
           \colvec[r]{2 \\ 2 \\ 0}
           -\proj{\colvec[r]{2 \\ 2 \\ 0}}{\spanof{\vec{\kappa}_1}}
           =\colvec[r]{2 \\ 2 \\ 0}
           -\frac{\colvec[r]{2 \\ 2 \\ 0}\dotprod\colvec[r]{1 \\ 5 \\ -1}}{%
                  \colvec[r]{1 \\ 5 \\ -1}\dotprod\colvec[r]{1 \\ 5 \\ -1}}
            \cdot\colvec[r]{1 \\ 5 \\ -1}                               \\
           &\hbox{}\quad=
           \colvec[r]{2 \\ 2 \\ 0}
           -\frac{12}{27}\cdot\colvec[r]{1 \\ 5 \\ -1}
           =\colvec[r]{14/9 \\ -2/9 \\ 4/9}                        
       \end{align*}
       Berikut anggota terakhirnya.
       \begin{align*}
        \vec{\kappa}_3    
           &=  
           \colvec[r]{1 \\ 0 \\ 0}
             -\proj{\colvec[r]{1 \\ 0 \\ 0}}{\spanof{\vec{\kappa}_1}}
             -\proj{\colvec[r]{1 \\ 0 \\ 0}}{\spanof{\vec{\kappa}_2}}       \\
           &\hbox{}\quad=
           \colvec[r]{1 \\ 0 \\ 0}
           -\frac{\colvec[r]{1 \\ 0 \\ 0}\dotprod\colvec[r]{1 \\ 5 \\ -1}}{%
                    \colvec[r]{1 \\ 5 \\ -1}\dotprod\colvec[r]{1 \\ 5 \\ -1}}
             \cdot\colvec[r]{1 \\ 5 \\ -1}
           -\frac{\colvec[r]{1 \\ 0 \\ 0}\dotprod\colvec[r]{14/9 \\ -2/9 \\  4/9}}{%
            \colvec[r]{14/9 \\ -2/9 \\ 4/9}\dotprod\colvec[r]{14/9 \\ -2/9 \\ 4/9}}
            \cdot\colvec[r]{14/9 \\ -2/9 \\ 4/9}              \\
         &=
         \colvec[r]{1 \\ 0 \\ 0}
           -\frac{1}{27}\cdot\colvec[r]{1 \\ 5 \\ -1}
           -\frac{7}{12}\cdot\colvec[r]{14/9 \\ -2/9 \\ 4/9}
         =\colvec[r]{1/18 \\ -1/18  \\ -4/18}
       \end{align*}
       Hasil $\vec{\kappa}_3$ ortogonal terhadap $\vec{\kappa}_1$ dan
       $\vec{\kappa}_2$.
       Karena itu, vektor tersebut ortogonal terhadap setiap vektor dalam
       rentang himpunan $\set{\vec{\kappa}_1,\vec{\kappa}_2}$, termasuk kedua
       vektor yang diberikan dalam soal.
     \end{answer}
  \recommended \item \label{exer:OrthoRepEasy} 
    Salah satu keunggulan basis ortogonal ialah bahwa basis tersebut
    menyederhanakan pencarian representasi suatu vektor terhadap basis itu.
    \begin{exparts}
      \partsitem Untuk vektor dan basis tak ortogonal bagi $\Re^2$ berikut,
        \begin{equation*}
          \vec{v}=\colvec[r]{2 \\ 3}
          \qquad
          B=\sequence{\colvec[r]{1 \\ 1},
                    \colvec[r]{1 \\ 0}}
        \end{equation*}
        mula-mula representasikan vektor terhadap basis tersebut.
        Kemudian proyeksikan vektor itu ke rentang setiap vektor basis,
        $\spanof{\smash{\vec{\beta}_1}}$ dan
        $\spanof{\smash{\vec{\beta}_2}}$.
      \partsitem Dengan basis ortogonal bagi $\Re^2$ berikut,
        \begin{equation*}
          K=\sequence{\colvec[r]{1 \\ 1},
                    \colvec[r]{1 \\ -1}}
        \end{equation*}
        representasikan vektor $\vec{v}$ yang sama terhadap basis tersebut.
        Kemudian proyeksikan vektor itu ke rentang setiap vektor basis.
        Perhatikan bahwa koefisien dalam representasi dan proyeksi sama.
      \partsitem Misalkan
        $K=\sequence{\vec{\kappa}_1,\ldots,\vec{\kappa}_k}$
        merupakan basis ortogonal bagi suatu subruang $\Re^n$.
        Buktikan bahwa bagi sebarang $\vec{v}$ dalam subruang tersebut,
        komponen ke-$i$ dari representasi $\rep{\vec{v}\,}{K}$ adalah
        koefisien skalar $(\vec{v}\dotprod\vec{\kappa}_i)/
               (\vec{\kappa}_i\dotprod\vec{\kappa}_i)$
        dari $\proj{\vec{v}\,}{\spanof{\vec{\kappa}_i}}$.
      \partsitem Buktikan bahwa
        $\vec{v}=\proj{\vec{v}\,}{\spanof{\vec{\kappa}_1}}
                   +\dots+\proj{\vec{v}\,}{\spanof{\vec{\kappa}_k}}$.
    \end{exparts}
    \begin{answer}
      \begin{exparts}
        \partsitem Representasinya dapat kita tentukan langsung.
          \begin{equation*}
            \colvec[r]{2 \\ 3}=3\cdot\colvec[r]{1 \\ 1}+(-1)\cdot\colvec[r]{1 \\ 0}
            \qquad
            \rep{\vec{v}\,}{B}=\colvec[r]{3 \\ -1}_B
          \end{equation*}
          Kedua proyeksinya juga mudah.
          \begin{align*}
            \proj{\colvec[r]{2 \\ 3}}{\spanof{\vec{\beta}_1}}
            &=\frac{\colvec[r]{2 \\ 3}\dotprod\colvec[r]{1 \\ 1}}{%
                   \colvec[r]{1 \\ 1}\dotprod\colvec[r]{1 \\ 1}}
              \cdot\colvec[r]{1 \\ 1}
            =\frac{5}{2}\cdot\colvec[r]{1 \\ 1}            \\
            \proj{\colvec[r]{2 \\ 3}}{\spanof{\vec{\beta}_2}}
            &=\frac{\colvec[r]{2 \\ 3}\dotprod\colvec[r]{1 \\ 0}}{%
                   \colvec[r]{1 \\ 0}\dotprod\colvec[r]{1 \\ 0}}
              \cdot\colvec[r]{1 \\ 0}
            =\frac{2}{1}\cdot\colvec[r]{1 \\ 0}
          \end{align*}
        \partsitem Seperti di atas, representasinya dapat kita tentukan langsung,
          \begin{equation*}
            \colvec[r]{2 \\ 3}=(5/2)\cdot\colvec[r]{1 \\ 1}
                             +(-1/2)\cdot\colvec[r]{1 \\ -1}
          \end{equation*}
          dan kedua proyeksinya mudah.
          \begin{align*}
            \proj{\colvec[r]{2 \\ 3}}{\spanof{\vec{\kappa}_1}}
            &=\frac{\colvec[r]{2 \\ 3}\dotprod\colvec[r]{1 \\ 1}}{%
                   \colvec[r]{1 \\ 1}\dotprod\colvec[r]{1 \\ 1}}
              \cdot\colvec[r]{1 \\ 1}
            =\frac{5}{2}\cdot\colvec[r]{1 \\ 1}                    \\
            \proj{\colvec[r]{2 \\ 3}}{\spanof{\vec{\kappa}_2}}
            &=\frac{\colvec[r]{2 \\ 3}\dotprod\colvec[r]{1 \\ -1}}{%
                   \colvec[r]{1 \\ -1}\dotprod\colvec[r]{1 \\ -1}}
              \cdot\colvec[r]{1 \\ -1}
            =\frac{-1}{2}\cdot\colvec[r]{1 \\ -1}
          \end{align*}
          Perhatikan kemunculan kembali $5/2$ dan $-1/2$.
        \partsitem Representasikan $\vec{v}$ terhadap basis
          \begin{equation*}
            \rep{\vec{v}\,}{K}=\colvec{r_1 \\ \vdots \\ r_k}
          \end{equation*}
          sehingga
          $\vec{v}=r_1\vec{\kappa}_1+\dots+r_k\vec{\kappa}_k$.
          Untuk menentukan $r_i$, ambil hasil kali titik kedua ruas dengan
          $\vec{\kappa}_i$.
          \begin{equation*}
            \vec{v}\dotprod\vec{\kappa}_i
               =\left(r_1\vec{\kappa}_1+\dots+r_k\vec{\kappa}_k\right)
                 \dotprod\vec{\kappa}_i
               =r_1\cdot 0+\dots+r_i\cdot(\vec{\kappa}_i\dotprod\vec{\kappa}_i)
                    +\dots+r_k\cdot 0
          \end{equation*}
          Menyelesaikannya untuk $r_i$ menghasilkan koefisien yang diinginkan.
        \partsitem Ini merupakan pernyataan ulang bagian sebelumnya.
      \end{exparts}
    \end{answer}
  \item 
     \textit{Ketaksamaan Bessel}.
     Tinjau himpunan-himpunan ortonormal berikut,
     \begin{equation*}
          B_1=\set{\vec{e}_1}
          \quad
          B_2=\set{\vec{e}_1,\vec{e}_2}
          \quad
          B_3=\set{\vec{e}_1,\vec{e}_2,\vec{e}_3}
          \quad
          B_4=\set{\vec{e}_1,\vec{e}_2,\vec{e}_3,\vec{e}_4}
     \end{equation*}
     beserta vektor $\vec{v}\in\Re^4$ yang komponennya $4$, $3$, $2$, dan $1$.
    \begin{exparts}
      \partsitem 
        Temukan koefisien $c_1$ bagi proyeksi $\vec{v}$ ke rentang vektor
        dalam $B_1$.
        Periksa bahwa $\absval{\vec{v}\,}^2\geq \absval{c_1}^2$.
      \partsitem 
        Temukan koefisien $c_1$ dan $c_2$ bagi proyeksi $\vec{v}$ ke rentang
        kedua vektor dalam $B_2$.
        Periksa bahwa $\absval{\vec{v}\,}^2\geq \absval{c_1}^2+\absval{c_2}^2$.
      \partsitem Temukan $c_1$, $c_2$, dan $c_3$ yang terkait dengan
        vektor-vektor dalam $B_3$, serta $c_1$, $c_2$, $c_3$, dan $c_4$ bagi
        vektor-vektor dalam $B_4$.
        Periksa bahwa
        $\absval{\vec{v}\,}^2\geq \absval{c_1}^2+\dots+\absval{c_3}^2$ dan
        bahwa
        $\absval{\vec{v}\,}^2\geq \absval{c_1}^2+\dots+\absval{c_4}^2$.
    \end{exparts}
    Tunjukkan bahwa ini berlaku secara umum:~jika
    $\set{\vec{\kappa}_1,\ldots,\vec{\kappa}_k}$
    merupakan himpunan ortonormal dan $c_i$ adalah koefisien proyeksi suatu
    vektor $\vec{v}$ dari ruang tersebut, maka
    $\absval{\vec{v}\,}^2\geq \absval{c_1}^2+\dots+\absval{c_k}^2$.
    \textit{Petunjuk}. Salah satu caranya ialah meninjau ketaksamaan
      $0\leq\absval{\vec{v}-(c_1\vec{\kappa}_1+\dots+c_k\vec{\kappa}_k)}^2$
      lalu menguraikan suku-suku $c$.
    \begin{answer}
      Pertama, $\absval{\vec{v}\,}^2=4^2+3^2+2^2+1^2=50$.
      \begin{exparts*}
        \partsitem $c_1=4$
        \partsitem $c_1=4$, $c_2=3$
        \partsitem $c_1=4$, $c_2=3$, $c_3=2$, $c_4=1$
      \end{exparts*}
      Untuk pembuktiannya, kita hanya akan mengerjakan kasus $k=2$ karena kasus
      yang sepenuhnya umum lebih rumit tetapi tidak lebih mencerahkan.
      Kita mengikuti petunjuk
      (ingat bahwa bagi sebarang vektor $\vec{w}$ berlaku
      $\absval{\vec{w}\,}^2=\vec{w}\dotprod\vec{w}$).
      \begin{align*}
        0
         &\leq
         \left(\vec{v}-\bigl(\frac{\vec{v}\dotprod\vec{\kappa}_1}{%
                              \vec{\kappa}_1\dotprod\vec{\kappa}_1}
                          \cdot\vec{\kappa}_1
                         +\frac{\vec{v}\dotprod\vec{\kappa}_2}{%
                              \vec{\kappa}_2\dotprod\vec{\kappa}_2}
                          \cdot\vec{\kappa}_2
                       \bigr)
         \right)
         \dotprod
         \left(\vec{v}-\bigl(\frac{\vec{v}\dotprod\vec{\kappa}_1}{%
                              \vec{\kappa}_1\dotprod\vec{\kappa}_1}
                          \cdot\vec{\kappa}_1
                         +\frac{\vec{v}\dotprod\vec{\kappa}_2}{%
                              \vec{\kappa}_2\dotprod\vec{\kappa}_2}
                          \cdot\vec{\kappa}_2
                       \bigr)
         \right)                                                       \\
         &=\begin{aligned}[t]
            &\vec{v}\dotprod\vec{v}
             -2\cdot\vec{v}
              \dotprod\left(\frac{\vec{v}\dotprod\vec{\kappa}_1}{%
                              \vec{\kappa}_1\dotprod\vec{\kappa}_1}
                          \cdot\vec{\kappa}_1
                       +\frac{\vec{v}\dotprod\vec{\kappa}_2}{%
                              \vec{\kappa}_2\dotprod\vec{\kappa}_2}
                          \cdot\vec{\kappa}_2\right)               \\
            &+\left(\frac{\vec{v}\dotprod\vec{\kappa}_1}{%
                              \vec{\kappa}_1\dotprod\vec{\kappa}_1}
                          \cdot\vec{\kappa}_1
                       +\frac{\vec{v}\dotprod\vec{\kappa}_2}{%
                              \vec{\kappa}_2\dotprod\vec{\kappa}_2}
                          \cdot\vec{\kappa}_2\right)
           \dotprod
           \left(\frac{\vec{v}\dotprod\vec{\kappa}_1}{%
                              \vec{\kappa}_1\dotprod\vec{\kappa}_1}
                          \cdot\vec{\kappa}_1
                       +\frac{\vec{v}\dotprod\vec{\kappa}_2}{%
                              \vec{\kappa}_2\dotprod\vec{\kappa}_2}
                          \cdot\vec{\kappa}_2\right)               
          \end{aligned}                                              \\
         &=
          \vec{v}\dotprod\vec{v}
          -2\cdot
             \left(\frac{\vec{v}\dotprod\vec{\kappa}_1}{%
                              \vec{\kappa}_1\dotprod\vec{\kappa}_1}
                          \cdot(\vec{v}\dotprod\vec{\kappa}_1)
                       +\frac{\vec{v}\dotprod\vec{\kappa}_2}{%
                              \vec{\kappa}_2\dotprod\vec{\kappa}_2}
                          \cdot(\vec{v}\dotprod\vec{\kappa}_2)\right)  \\
          &\quad+\left((\frac{\vec{v}\dotprod\vec{\kappa}_1}{%
                              \vec{\kappa}_1\dotprod\vec{\kappa}_1})^2
                          \cdot(\vec{\kappa}_1\dotprod\vec{\kappa}_1)
                       +(\frac{\vec{v}\dotprod\vec{\kappa}_2}{%
                              \vec{\kappa}_2\dotprod\vec{\kappa}_2})^2
                          \cdot(\vec{\kappa}_2\dotprod\vec{\kappa}_2)\right)
      \end{align*}
      (Kedua suku silang pada bagian ketiga dari baris ketiga bernilai nol
      karena $\vec{\kappa}_1$ dan $\vec{\kappa}_2$ ortogonal.)
      Hasilnya kini diperoleh dengan menggabungkan suku-suku sejenis dan
      memperhatikan bahwa
      $\vec{\kappa}_1\dotprod\vec{\kappa}_1=1$ dan
      $\vec{\kappa}_2\dotprod\vec{\kappa}_2=1$ karena vektor-vektor ini
      merupakan anggota suatu himpunan ortonormal.
    \end{answer}
  \item
    Buktikan atau sangkal:~setiap vektor dalam \( \Re^n \) termuat dalam suatu
    basis ortogonal.
    \begin{answer}
      Pernyataan ini benar, kecuali bagi vektor nol.
      Setiap vektor dalam \( \Re^n \), selain vektor nol, termuat dalam suatu
      basis, dan basis tersebut dapat diortogonalisasi.
     \end{answer}
  \item 
    Tunjukkan bahwa kolom-kolom suatu matriks \( \nbyn{n} \) membentuk himpunan
    ortonormal jika dan hanya jika invers matriks tersebut adalah transposnya.
    Berikan sebuah matriks semacam itu.
    \begin{answer} 
      Kasus $\nbyn{3}$ memberikan gagasannya.
      Himpunan
      \begin{equation*}
        \set{\colvec{a \\ d \\ g},\colvec{b \\ e \\ h},\colvec{c \\ f \\ i}}
      \end{equation*}
      ortonormal jika dan hanya jika kesembilan syarat berikut semuanya berlaku,
      \begin{equation*}
        \begin{array}{ccc}
          \rowvec{a &d &g}\dotprod\colvec{a \\ d \\ g}=1
            &\rowvec{a &d &g}\dotprod\colvec{b \\ e \\ h}=0
            &\rowvec{a &d &g}\dotprod\colvec{c \\ f \\ i}=0    \\
          \rowvec{b &e &h}\dotprod\colvec{a \\ d \\ g}=0
            &\rowvec{b &e &h}\dotprod\colvec{b \\ e \\ h}=1
            &\rowvec{b &e &h}\dotprod\colvec{c \\ f \\ i}=0    \\
          \rowvec{c &f &i}\dotprod\colvec{a \\ d \\ g}=0
            &\rowvec{c &f &i}\dotprod\colvec{b \\ e \\ h}=0
            &\rowvec{c &f &i}\dotprod\colvec{c \\ f \\ i}=1    
        \end{array}
      \end{equation*}
      (ketiga syarat di kiri bawah berlebih, tetapi tetap benar).
      Syarat-syarat tersebut, pada gilirannya, berlaku jika dan hanya jika
      \begin{equation*}
        \begin{mat}
          a &d &g \\
          b &e &h \\
          c &f &i
        \end{mat}
        \begin{mat}
          a &b &c  \\
          d &e &f  \\
          g &h &i
        \end{mat}
        =
        \begin{mat}
          1 &0 &0  \\
          0 &1 &0  \\
          0 &0 &1
        \end{mat}
      \end{equation*}
      sebagaimana yang disyaratkan.

      Berikut sebuah contoh; invers matriks ini adalah transposnya.
      \begin{equation*}
        \begin{mat}[r]
          1/\sqrt{2}  &1/\sqrt{2}  &0  \\
         -1/\sqrt{2}  &1/\sqrt{2}  &0  \\
          0           &0           &1
        \end{mat}
      \end{equation*}
    \end{answer}
  \item 
    Apakah bukti \nearbytheorem{th:OrthoIsInd} gagal mempertimbangkan
    kemungkinan bahwa himpunan vektornya kosong (yakni, $k=0$)?
    \begin{answer}
      Jika himpunannya kosong, maka penjumlahan di ruas kiri merupakan
      kombinasi linear himpunan vektor kosong, yang menurut definisi berjumlah
      vektor nol.
      Dalam kalimat kedua, tidak ada $i$ seperti itu, sehingga implikasi
      `jika \ldots maka \ldots' benar secara hampa.
    \end{answer}
  \item 
    \nearbytheorem{th:GramSchmidt} mendeskripsikan perubahan basis dari
    sebarang basis \( B=\sequence{\vec{\beta}_1,\ldots,\vec{\beta}_k} \)
    menjadi basis ortogonal
    \( K=\sequence{\vec{\kappa}_1,\ldots,\vec{\kappa}_k} \).
    Tinjau matriks perubahan basis $\rep{\identity}{B,K}$.
    \begin{exparts}
      \partsitem Buktikan bahwa matriks $\rep{\identity}{K,B}$, yang mengubah
        basis dalam arah yang berlawanan dengan arah dalam teorema, berbentuk
        segitiga atas\Dash semua entri di bawah diagonal utamanya nol.
      \partsitem Buktikan bahwa invers matriks segitiga atas juga segitiga atas
        (tentu saja jika matriks tersebut dapat diinverskan).
        Ini menunjukkan bahwa matriks $\rep{\identity}{B,K}$, yang mengubah
        basis dalam arah yang dideskripsikan teorema, berbentuk segitiga atas.
    \end{exparts}
    \begin{answer} 
      \begin{exparts}
        \partsitem Sebagian argumen induksi yang membuktikan
          \nearbytheorem{th:GramSchmidt} memeriksa bahwa $\vec{\kappa}_i$
          berada dalam rentang
          $\sequence{\vec{\beta}_1,\dots,\vec{\beta}_i}$.
          (Kasus $i=3$ dalam bukti menjadi ilustrasi.)
          Jadi, dalam matriks perubahan basis $\rep{\identity}{K,B}$, kolom
          ke-$i$, yakni $\rep{\vec{\kappa}_i}{B}$, memiliki komponen
          $i+1$ hingga $k$ yang bernilai nol.
        \partsitem Salah satu cara melihatnya ialah mengingat kembali prosedur
          komputasi yang digunakan untuk mencari invers.
          Kita menuliskan matriks, menuliskan matriks identitas di sebelahnya,
          lalu melakukan reduksi Gauss--Jordan.
          Jika matriks pada awalnya berbentuk segitiga atas, reduksi
          Gauss--Jordan hanya melibatkan bagian Jordan, dan langkah-langkah ini,
          ketika diterapkan pada matriks identitas, akan menghasilkan matriks
          invers segitiga atas.
      \end{exparts}
    \end{answer}
 \item \label{exer:GramSchmidt} 
    Lengkapi argumen induksi dalam bukti \nearbytheorem{th:GramSchmidt}.
    \begin{answer}
      Untuk langkah induksi, kita mengasumsikan bahwa bagi semua $j$ dalam
      $[1..i]$, ketiga syarat berikut berlaku bagi setiap $\vec{\kappa}_j$:
      (i)~setiap $\vec{\kappa}_j$ tak nol,
      (ii)~setiap $\vec{\kappa}_j$ merupakan kombinasi linear vektor-vektor
         $\vec{\beta}_1,\dots,\vec{\beta}_j$, dan
      (iii)~setiap $\vec{\kappa}_j$ ortogonal terhadap semua
         $\vec{\kappa}_m$ sebelumnya (yakni, dengan $m<j$).
      Dengan hipotesis-hipotesis induksi tersebut, tinjau
      $\vec{\kappa}_{i+1}$.
      \begin{align*}
        \vec{\kappa}_{i+1}
         &=
        \vec{\beta}_{i+1}
          -\proj{\vec{\beta}_{i+1}}{\spanof{\vec{\kappa}_1}}
          -\proj{\vec{\beta}_{i+1}}{\spanof{\vec{\kappa}_2}}
          -\dots
          -\proj{\vec{\beta}_{i+1}}{\spanof{\vec{\kappa}_i}}  \\
         &=
        \vec{\beta}_{i+1}
          -\frac{\vec{\beta}_{i+1}\dotprod\vec{\kappa}_1}{%
                  \vec{\kappa}_1\dotprod\vec{\kappa}_1}
              \cdot\vec{\kappa}_1
          -\frac{\vec{\beta}_{i+1}\dotprod\vec{\kappa}_2}{%
                  \vec{\kappa}_2\dotprod\vec{\kappa}_2}
              \cdot\vec{\kappa}_2
          -\dots
          -\frac{\vec{\beta}_{i+1}\dotprod\vec{\kappa}_i}{%
                  \vec{\kappa}_i\dotprod\vec{\kappa}_i}
              \cdot\vec{\kappa}_i
       \end{align*}
       Berdasarkan asumsi induksi~(ii), kita dapat menguraikan setiap
       $\vec{\kappa}_j$ menjadi kombinasi linear
       $\vec{\beta}_1,\ldots,\vec{\beta}_j$,
       \begin{align*}
         &=
        \vec{\beta}_{i+1}
          -\frac{\vec{\beta}_{i+1}\dotprod\vec{\kappa}_1}{%
                  \vec{\kappa}_1\dotprod\vec{\kappa}_1}
              \cdot\vec{\beta}_1                          \\
          &\quad-\frac{\vec{\beta}_{i+1}\dotprod\vec{\kappa}_2}{%
                  \vec{\kappa}_2\dotprod\vec{\kappa}_2}
              \cdot
                 \left(\text{
                    kombinasi linear $\vec{\beta}_1,\vec{\beta}_2$}
                 \right)                                                 \\
          &\quad-\;\cdots\;
          -\frac{\vec{\beta}_{i+1}\dotprod\vec{\kappa}_i}{%
                  \vec{\kappa}_i\dotprod\vec{\kappa}_i}
              \cdot
                 \left(\text{
                    kombinasi linear $\vec{\beta}_1,\dots,\vec{\beta}_i$}
                 \right)
       \end{align*}
       Pecahan-pecahan tersebut adalah skalar, sehingga ini merupakan kombinasi
       linear dari kombinasi-kombinasi linear
       $\vec{\beta}_1,\dots,\vec{\beta}_{i+1}$.
       Jadi, ini hanyalah kombinasi linear
       $\vec{\beta}_1,\dots,\vec{\beta}_{i+1}$.
       Sekarang, (i)~hasilnya tidak mungkin vektor nol, karena persamaan itu
       akan mendeskripsikan hubungan linear tak trivial di antara
       vektor-vektor $\vec{\beta}$ yang diberikan sebagai anggota suatu basis
       (hubungannya tak trivial karena koefisien $\vec{\beta}_{i+1}$ adalah
       $1$).
       Selain itu, (ii)~persamaan tersebut memberikan $\vec{\kappa}_{i+1}$
       sebagai kombinasi $\vec{\beta}_1,\dots,\vec{\beta}_{i+1}$.
       Akhirnya, untuk~(iii), tinjau
       $\vec{\kappa}_j\dotprod\vec{\kappa}_{i+1}$; seperti dalam kasus $i=3$,
       hasil kali titik $\vec{\kappa}_j$ dengan
       $\vec{\kappa}_{i+1}=\vec{\beta}_{i+1}
          -\proj{\vec{\beta}_{i+1}}{\spanof{\vec{\kappa}_1}}
          -\dots
          -\proj{\vec{\beta}_{i+1}}{\spanof{\vec{\kappa}_i}}$
         dapat ditulis ulang sehingga memberikan dua jenis suku,
       $\vec{\kappa}_j\dotprod
           \left(\vec{\beta}_{i+1}
                  -\proj{\vec{\beta}_{i+1}}{\spanof{\vec{\kappa}_j}}
           \right)$
       (yang bernilai nol karena proyeksinya ortogonal), dan
       $\vec{\kappa}_j\dotprod
               \proj{\vec{\beta}_{i+1}}{\spanof{\vec{\kappa}_m}}$
       dengan $m\neq j$ dan $m<i+1$
       (yang bernilai nol karena berdasarkan hipotesis~(iii), vektor
       $\vec{\kappa}_j$ dan $\vec{\kappa}_m$ ortogonal).
     \end{answer}
\index{proses Gram--Schmidt|)}
\end{exercises}



















\subsectionoptional{Proyeksi ke Sebuah Subruang}
\noindent\textit{Subbagian ini menggunakan materi dari subbagian opsional
     sebelumnya tentang Menggabungkan Subruang.}

Subbagian-subbagian sebelumnya memproyeksikan suatu vektor ke sebuah garis
dengan menguraikannya menjadi dua bagian:~bagian pada garis,
$\proj{\vec{v}\,}{\spanof{\vec{s}\,}}$, dan sisanya,
$\vec{v}-\proj{\vec{v}\,}{\spanof{\vec{s}\,}}$.
Untuk memperumum proyeksi ke subruang sebarang, kita akan mengikuti gagasan
penguraian ini.

\begin{definition}
\label{def:ProjIntoMAlongN}
Misalkan suatu ruang vektor merupakan jumlah langsung \( V=M\directsum N \).
Untuk sebarang \( \vec{v}\in V \) dengan $\vec{v}=\vec{m}+\vec{n}$, di mana
\( \vec{m}\in M,\,\vec{n}\in N \),
\definend{proyeksi \( \vec{v} \) ke \( M \) sepanjang \( N \)}%
\index{proyeksi!ke sebuah subruang}\index{proyeksi!sepanjang sebuah subruang}
adalah
$\proj{\vec{v}\,}{M,N}=\vec{m}$.
\end{definition}

Definisi ini berlaku pada ruang-ruang yang belum memiliki definisi ortogonal
yang siap digunakan.
(Definisi keortogonalan bagi ruang selain $\Re^n$ sangat mungkin dibuat,
tetapi kita belum menjumpainya dalam buku ini.)

\begin{example} \label{ex:OrthProjOne}
Ruang \( \matspace_{\nbyn{2}} \) yang terdiri atas matriks-matriks
$\nbyn{2}$ merupakan jumlah langsung kedua subruang berikut.
\begin{equation*}
   M=\set{\begin{mat}
              a  &b  \\
              0  &0
            \end{mat} \suchthat a,b\in\Re }
  \qquad
   N=\set{\begin{mat}
              0  &0  \\
              c  &d
            \end{mat} \suchthat c,d\in\Re }
\end{equation*}
Untuk memproyeksikan
\begin{equation*}
  A=\begin{mat}[r]
          3  &1  \\
          0  &4
  \end{mat}
\end{equation*}
ke $M$ sepanjang $N$, pertama-tama kita tetapkan basis bagi kedua subruang.
\begin{equation*}
  B_{M}=\sequence{
    \begin{mat}[r]
      1  &0  \\
      0  &0      
    \end{mat},
    \begin{mat}[r]
      0  &1  \\
      0  &0
    \end{mat}
                   }
  \qquad
  B_{N}=\sequence{
    \begin{mat}[r]
      0  &0  \\
      1  &0      
    \end{mat},
    \begin{mat}[r]
      0  &0  \\
      0  &1
    \end{mat}
                   }
\end{equation*}
Konkatenasi keduanya,
\begin{equation*}
  B=\cat{B_{M}}{B_{N}}=\sequence{
    \begin{mat}[r]
      1  &0  \\
      0  &0      
    \end{mat},
    \begin{mat}[r]
      0  &1  \\
      0  &0
    \end{mat},
    \begin{mat}[r]
      0  &0  \\
      1  &0      
    \end{mat},
    \begin{mat}[r]
      0  &0  \\
      0  &1
    \end{mat}
                   }
\end{equation*}
merupakan basis bagi seluruh ruang karena \( \matspace_{\nbyn{2}} \)
merupakan jumlah langsung tersebut.
Jadi, kita dapat menggunakannya untuk merepresentasikan $A$.
\begin{equation*}
  \begin{mat}[r]
          3  &1  \\
          0  &4
  \end{mat}=
    3\cdot\begin{mat}[r]
      1  &0  \\
      0  &0      
    \end{mat}
    +1\cdot\begin{mat}[r]
      0  &1  \\
      0  &0
    \end{mat}
    +0\cdot\begin{mat}[r]
      0  &0  \\
      1  &0      
    \end{mat}
    +4\cdot\begin{mat}[r]
      0  &0  \\
      0  &1
    \end{mat}
\end{equation*}
Proyeksi $A$ ke $M$ sepanjang $N$ mempertahankan bagian $M$ dan membuang
bagian $N$.
\begin{equation*}
  \proj{\begin{mat}[r]
          3  &1  \\
          0  &4
  \end{mat} }{M,N}
  =
    3\cdot\begin{mat}[r]
      1  &0  \\
      0  &0      
    \end{mat}
    +1\cdot\begin{mat}[r]
      0  &1  \\
      0  &0
    \end{mat}
    =\begin{mat}[r]
      3  &1  \\
      0  &0      
    \end{mat}
\end{equation*}
\end{example}

\begin{example}  \label{ex:ProjIntoMAlongNGener}
Kedua subskrip pada $\proj{\vec{v}\,}{M,N}$ sama-sama penting.
Subskrip pertama, $M$, penting karena hasil proyeksi merupakan anggota $M$.
Sebagai contoh yang menunjukkan pentingnya subskrip kedua, tetapkan subruang
bidang dari $\Re^3$ berikut beserta basisnya.
\begin{equation*}
  M=\set{\colvec{x \\ y \\ z}\suchthat y-2z=0}
  \qquad
  B_M=\sequence{
                \colvec[r]{1 \\ 0 \\ 0},
                \colvec[r]{0 \\ 2 \\ 1} }
\end{equation*}
Kita akan membandingkan proyeksi elemen~$\Re^3$ berikut
\begin{equation*}
  \vec{v}=\colvec[r]{2 \\ 2 \\ 5}
\end{equation*}
ke $M$ sepanjang kedua subruang berikut
(verifikasi bahwa \( \Re^3=M\directsum N \) dan
\( \Re^3=M\directsum \hat{N} \) bersifat rutin).
\begin{equation*}
  N=\set{k\colvec[r]{0 \\ 0 \\ 1}\suchthat k\in\Re}
  \qquad
  \hat{N}=\set{k\colvec[r]{0 \\ 1 \\ -2}\suchthat k\in\Re}
\end{equation*}
Berikut basis alami bagi $N$ dan $\hat{N}$.
\begin{equation*}
  B_N=\sequence{
                \colvec[r]{0 \\ 0 \\ 1} }
  \qquad
  B_{\hat{N}}=\sequence{
                \colvec[r]{0 \\ 1 \\ -2} }
\end{equation*}
Untuk memproyeksikan ke $M$ sepanjang~$N$, representasikan $\vec{v}$ terhadap
konkatenasi $\cat{B_M}{B_N}$,
\begin{equation*}
  \colvec[r]{2 \\ 2 \\ 5}=
  2\cdot\colvec[r]{1 \\ 0 \\ 0}+
  1\cdot\colvec[r]{0 \\ 2 \\ 1}+
  4\cdot\colvec[r]{0 \\ 0 \\ 1}
\end{equation*}
dan buang suku $N$.
\begin{equation*}
  \proj{\vec{v}\,}{M,N}
  =2\cdot\colvec[r]{1 \\ 0 \\ 0}+
  1\cdot\colvec[r]{0 \\ 2 \\ 1}
  =\colvec[r]{2 \\ 2 \\ 1}
\end{equation*}
Untuk memproyeksikan ke~$M$ sepanjang~$\hat{N}$, representasikan $\vec{v}$
terhadap $\cat{B_M}{B_{\hat{N}}}$,
\begin{equation*}
  \colvec[r]{2 \\ 2 \\ 5}=
  2\cdot\colvec[r]{1 \\ 0 \\ 0}
  +(9/5)\cdot\colvec[r]{0 \\ 2 \\ 1}
  -(8/5)\cdot\colvec[r]{0 \\ 1 \\ -2}
\end{equation*}
dan hilangkan bagian $\hat{N}$.
\begin{equation*}
  \proj{\vec{v}\,}{M,\hat{N}}
  =
  2\cdot\colvec[r]{1 \\ 0 \\ 0}
  +(9/5)\cdot\colvec[r]{0 \\ 2 \\ 1}
  =\colvec[r]{2 \\ 18/5 \\ 9/5}
\end{equation*}
Jadi, memproyeksikan sepanjang subruang yang berbeda dapat memberikan hasil
yang berbeda.

Gambar-gambar berikut membandingkan kedua pemetaan.
Keduanya menunjukkan bahwa proyeksi tersebut benar-benar `ke' bidang dan
`sepanjang' garis.
\begin{center}  \small
  \begin{tabular}{@{}c@{}}\includegraphics{map/mp/ch3.38}\end{tabular}
  \qquad
  \begin{tabular}{@{}c@{}}\includegraphics{map/mp/ch3.39}\end{tabular}
\end{center}
Perhatikan bahwa proyeksi sepanjang $N$ tidak ortogonal karena ada anggota
bidang $M$ yang tidak ortogonal terhadap garis putus-putus.
Namun, proyeksi sepanjang $\hat{N}$ bersifat ortogonal.
\end{example}

Kita telah melihat dua operasi proyeksi: proyeksi ortogonal ke sebuah garis,
serta proyeksi dalam subbagian ini ke suatu~$M$ sepanjang suatu~$N$; wajar jika
kita bertanya apakah keduanya berkaitan.
Gambar sebelah kanan di atas menunjukkan jawabannya\Dash proyeksi ortogonal ke
sebuah garis merupakan kasus khusus proyeksi dalam subbagian ini; proyeksi itu
berlangsung sepanjang subruang yang tegak lurus terhadap garis.
\begin{center}  \small
  \includegraphics{map/mp/ch3.40}
\end{center}
% In addition to pointing out that projection along a subspace is a 
% generalization, this scheme shows how to define orthogonal projection into
% any subspace of $\Re^n$, of any dimension. 

\begin{definition} \label{def:OrthComp}
\definend{Komplemen ortogonal\/}\index{ortogonal!komplemen}%
\index{subruang komplementer!ortogonal}\index{perp, dari suatu subruang}
dari suatu subruang \( M \) dalam $\Re^n$ adalah
\begin{equation*}
  M^\perp=\set{\vec{v}\in\Re^n\suchthat
                 \vec{v} \text{ tegak lurus terhadap semua vektor dalam }M}
\end{equation*}
(dibaca ``\( M \) perp'').
\definend{Proyeksi ortogonal}\index{ortogonal!proyeksi}%
\index{proyeksi!ortogonal} \( \proj{\vec{v}\,}{M} \) dari suatu vektor adalah
proyeksinya ke $M$ sepanjang \( M^\perp \).
\end{definition}


\begin{example} \label{ex:OrthoCompOne}
Dalam \( \Re^3 \), untuk mencari komplemen ortogonal bidang
\begin{equation*}
   P=\set{\colvec{x \\ y \\ z}\suchthat 3x+2y-z=0 }
\end{equation*}
kita mulai dengan suatu basis bagi \( P \).
\begin{equation*}
   B=\sequence{\colvec[r]{1 \\ 0 \\ 3},
               \colvec[r]{0 \\ 1 \\ 2} }
\end{equation*}
Setiap \( \vec{v} \) yang tegak lurus terhadap semua vektor dalam $B$ juga
tegak lurus terhadap setiap vektor dalam rentang $B$
(buktinya menjadi \nearbyexercise{exer:PerpOnBasisImplPerp}).
Oleh karena itu, subruang $P^\perp$ terdiri atas vektor-vektor yang memenuhi
kedua syarat berikut.
\begin{equation*}
  \colvec[r]{1 \\ 0 \\ 3}\dotprod\colvec{v_1 \\ v_2 \\ v_3}=0
  \qquad
  \colvec[r]{0 \\ 1 \\ 2}\dotprod\colvec{v_1 \\ v_2 \\ v_3}=0
\end{equation*}
Syarat-syarat tersebut memberikan suatu sistem linear.
\begin{equation*}
  P^\perp
    =\set{\colvec{v_1 \\ v_2 \\ v_3}\suchthat \begin{mat}[r]
                                                      1  &0  &3  \\
                                                      0  &1  &2
                                                    \end{mat}
                                                    \colvec{v_1 \\ v_2 \\ v_3}
                                                    =\colvec[r]{0 \\ 0} }
\end{equation*}
Jadi, tinggal dicari ruang nol dari pemetaan yang direpresentasikan oleh
matriks tersebut, yakni menghitung himpunan solusi sistem linear homogen.
\begin{equation*}
  \begin{linsys}{3}
    v_1 &  &    &+ &3v_3 &= &0 \\
        &  &v_2 &+ &2v_3 &= &0
  \end{linsys}
  \quad\Longrightarrow\quad
  P^\perp=\set{k\colvec[r]{-3 \\ -2 \\ 1}\suchthat k\in\Re}
\end{equation*}
% Instead of the term orthogonal complement,
% this is sometimes called the line \definend{normal}\index{normal} %
% to the plane.
\end{example}

\begin{example} \label{ex:OrthoCompTwo}
Jika \( M \) adalah subruang bidang~\( xy \) dari \( \Re^3 \), apa bentuk
\( M^\perp \)?
Tanggapan awal yang lazim adalah bahwa \( M^\perp \) merupakan bidang~\( yz \),
tetapi itu tidak benar karena beberapa vektor dari bidang~\( yz \) tidak tegak
lurus terhadap setiap vektor dalam bidang~\( xy \).
\begin{center}  \small
    {\small $\colvec[r]{1 \\ 1 \\ 0} \not\perp \colvec[r]{0 \\ 3 \\ 2}$}
   \quad
  \begin{tabular}{@{}c@{}}\includegraphics{map/mp/ch3.41}\end{tabular}
   \quad
   {\small $\displaystyle 
       \theta=\arccos(\frac{1\cdot 0+1\cdot 3+0\cdot 2}{
                            \sqrt{2}\cdot\sqrt{13}})
    \approx\text{$0.94$~rad}$}
\end{center}
Sebaliknya, \( M^\perp \) adalah sumbu~\( z \), karena dengan mengikuti contoh
sebelumnya dan mengambil basis alami bagi bidang~$xy$ diperoleh
\begin{equation*}
  M^\perp
  =
  \set{\colvec{x \\ y \\ z}\suchthat \begin{mat}[r]
                                         1  &0  &0  \\
                                         0  &1  &0
                                       \end{mat}
                                       \colvec{x \\ y \\ z}=
                                       \colvec[r]{0 \\ 0} \;}
  =
  \set{\colvec{x \\ y \\ z}\suchthat x=0 \text{\ dan\ }y=0  }
\end{equation*}
\end{example}

% The two examples that we've seen since \nearbydefinition{def:OrthComp}
% illustrate the first sentence in that definition.
% The next result justifies the second sentence.

\begin{lemma} \label{le:OrthoProjWellDefd}
Jika $M$ merupakan subruang $\Re^n$, maka komplemen ortogonalnya $M^\perp$ juga
merupakan subruang.
Ruang tersebut merupakan jumlah langsung keduanya,
\( \Re^n=M\directsum M^\perp \).
Bagi sebarang $\vec{v}\in\Re^n$, vektor
$\vec{v}-\proj{\vec{v}\,}{M}$ tegak lurus terhadap setiap vektor dalam $M$.
\end{lemma}

\begin{proof}
Pertama, komplemen ortogonal $M^\perp$ merupakan subruang $\Re^n$:
vektor nol berada di dalamnya, dan jika $\vec{u},\vec{w}\in M^\perp$ serta
$a,b\in\Re$, maka bagi setiap $\vec{m}\in M$ berlaku
$(a\vec{u}+b\vec{w})\dotprod\vec{m}
=a(\vec{u}\dotprod\vec{m})+b(\vec{w}\dotprod\vec{m})=0$.
 
Untuk menunjukkan bahwa ruang~$\R^n$ merupakan jumlah langsung keduanya,
mulailah dengan sebarang basis
\( B_M=\sequence{\vec{\mu}_1,\dots,\vec{\mu}_k} \) bagi \( M \).
Perluas basis tersebut menjadi basis
%\( B=\sequence{\vec{\mu}_1,\dots,\vec{\mu}_k,
%   \vec{\nu}_1,\dots,\vec{\nu}_{n-k} } \)
seluruh ruang, lalu terapkan proses Gram--Schmidt untuk memperoleh basis
ortogonal \( K=\sequence{\vec{\kappa}_1,\dots,\vec{\kappa}_n} \) bagi
\( \Re^n \).
Basis $K$ ini merupakan konkatenasi dua basis:
\( \sequence{\vec{\kappa}_1,\dots,\vec{\kappa}_k} \)
yang beranggota sama banyak, yaitu~$k$, dengan $B_M$, dan
\( D=\sequence{\vec{\kappa}_{k+1},\dots,\vec{\kappa}_n} \).
Yang pertama merupakan basis bagi $M$. Jadi, jika kita menunjukkan bahwa yang
kedua merupakan basis bagi \( M^\perp \), maka seluruh ruang merupakan jumlah
langsung tersebut.

\nearbyexercise{exer:OrthoRepEasy} dari subbagian sebelumnya membuktikan hal
berikut tentang sebarang basis ortogonal:~setiap vektor $\vec{v}$ dalam ruang
merupakan jumlah proyeksi ortogonalnya ke garis-garis yang direntang oleh
vektor-vektor basis.
\begin{equation*}
  \vec{v}=\proj{\vec{v}\,}{\spanof{\vec{\kappa}_1}}
           +\dots+\proj{\vec{v}\,}{\spanof{\vec{\kappa}_n}}
\tag*{($*$)}\end{equation*}
Untuk memeriksanya, representasikan vektor sebagai
$\vec{v}=r_1\vec{\kappa}_1+\dots+r_n\vec{\kappa}_n$,
ambil hasil kali titik kedua ruas dengan $\vec{\kappa}_i$,
$
   \vec{v}\dotprod\vec{\kappa}_i
      =\left(r_1\vec{\kappa}_1+\dots+r_n\vec{\kappa}_n\right)
          \dotprod\vec{\kappa}_i
      =r_1\cdot 0+\dots+r_i\cdot(\vec{\kappa}_i\dotprod\vec{\kappa}_i)
            +\dots+r_n\cdot 0
$,
lalu selesaikan untuk memperoleh
$r_i=(\vec{v}\dotprod\vec{\kappa}_i)/(\vec{\kappa}_i\dotprod\vec{\kappa}_i)$,
sebagaimana yang diinginkan.

Setiap anggota rentang
\( D \) %\sequence{\vec{\kappa}_{k+1},\dots,\vec{\kappa}_n} \)
ortogonal terhadap sebarang vektor dalam $M$, sehingga rentang \( D \)
merupakan subhimpunan \( M^\perp \).
Untuk menunjukkan bahwa \( D \) merupakan basis bagi $M^\perp$, kita hanya
perlu menunjukkan inklusi sebaliknya, yaitu bahwa setiap
$\vec{w}\in M^\perp$ merupakan elemen rentang~\( D \).
Paragraf sebelumnya dapat digunakan untuk ini.
Setiap $\vec{w}\in M^\perp$ memberikan proyeksi-proyeksi berikut ke
vektor-vektor basis dari $M$:
$\proj{\vec{w}\,}{\spanof{\vec{\kappa}_1}}=\zero,
\dots,\,\proj{\vec{w}\,}{\spanof{\vec{\kappa}_k}}=\zero$.
Oleh karena itu, persamaan~($*$) memberikan bahwa $\vec{w}$ merupakan kombinasi
linear \( \vec{\kappa}_{k+1},\dots,\vec{\kappa}_n \).
Jadi, \( D \) merupakan basis bagi $M^\perp$ dan $\Re^n$ merupakan jumlah
langsung keduanya.

Kalimat terakhir lema dibuktikan dengan cara yang hampir sama.
Tuliskan $\vec{v}=\proj{\vec{v}\,}{\spanof{\vec{\kappa}_1}}
           +\dots+\proj{\vec{v}\,}{\spanof{\vec{\kappa}_n}}$.
Kemudian $\proj{\vec{v}\,}{M}$ hanya mempertahankan bagian $M$ dan membuang
bagian $M^\perp$:
$\proj{\vec{v}\,}{M}=\proj{\vec{v}\,}{\spanof{\vec{\kappa}_{1}}}
                       +\dots+\proj{\vec{v}\,}{\spanof{\vec{\kappa}_k}}$.
Oleh karena itu, $\vec{v}-\proj{\vec{v}\,}{M}$ merupakan kombinasi linear
elemen-elemen $M^\perp$, sehingga tegak lurus terhadap setiap vektor dalam $M$.
\end{proof}

Untuk suatu subruang, kita dapat menghitung proyeksi ortogonal ke subruang
tersebut dengan mengikuti langkah-langkah bukti tadi: mencari sebuah basis,
memperluasnya menjadi basis bagi seluruh ruang, menerapkan Gram--Schmidt untuk
memperoleh basis ortogonal, dan memproyeksikan ke setiap subruang linear.
Namun, sebagai gantinya kita akan menggunakan rumus yang praktis.

\begin{theorem} \label{th:OrthProjMat}
Misalkan \( M \) merupakan subruang \( \Re^n \) dengan basis
\( \sequence{\vec{\beta}_1,\dots,\vec{\beta}_k} \), dan misalkan \( A \)
adalah matriks yang kolom-kolomnya merupakan vektor-vektor \( \vec{\beta} \).
Bagi sebarang $\vec{v}\in\Re^n$, proyeksi ortogonalnya adalah
$\proj{\vec{v}\,}{M}=c_1\vec{\beta}_1+\dots+c_k\vec{\beta}_k$,
di mana koefisien-koefisien $c_i$ merupakan entri-entri vektor
$(\trans{A}A)^{-1}\trans{A}\cdot\vec{v}$.
Artinya,
$\proj{\vec{v}\,}{M}=A(\trans{A}A)^{-1}\trans{A}\cdot\vec{v}$.
\end{theorem}

\begin{proof}
Vektor \( \proj{\vec{v}}{M} \) merupakan anggota \( M \), sehingga merupakan
kombinasi linear vektor-vektor basis,
\( c_1\cdot\vec{\beta}_1+\dots+c_k\cdot\vec{\beta}_k \).
Karena kolom-kolom $A$ adalah vektor-vektor $\vec{\beta}$, terdapat
\( \vec{c}\in\Re^k \) sedemikian sehingga
\( \proj{\vec{v}\,}{M}=A\vec{c} \).
% (this is expressed compactly with matrix
% multiplication as in
% \nearbyexample{ex:OrthoCompOne} and~\ref{ex:OrthoCompTwo}).
Untuk mencari $\vec{c}$, perhatikan bahwa vektor
\( \vec{v}-\proj{\vec{v}\,}{M} \) tegak lurus terhadap setiap anggota basis,
sehingga
 % (again, expressed compactly).
\begin{equation*}
  \zero
  =\trans{A}\bigl( \vec{v}-A\vec{c} \bigr) 
  =\trans{A}\vec{v}-\trans{A}A\vec{c}
\end{equation*}
dan menyelesaikannya memberikan berikut
(menunjukkan bahwa \( \trans{A}A \) dapat diinverskan menjadi sebuah latihan).
\begin{equation*}
  \vec{c}=\bigl( \trans{A}A\bigr)^{-1}\trans{A}\cdot\vec{v}
\end{equation*}
Oleh karena itu,
$\proj{\vec{v}\,}{M}=A\cdot\vec{c}=A(\trans{A}A)^{-1}\trans{A}\cdot\vec{v}$,
sebagaimana yang disyaratkan.
\end{proof}

\begin{example}
Untuk memproyeksikan secara ortogonal vektor ini ke subruang berikut,
\begin{equation*}
  \vec{v}=\colvec[r]{1 \\ -1 \\ 1}
  \qquad
  P=\set{\colvec{x \\ y \\ z}\suchthat x+z=0}
\end{equation*}
mula-mula susun matriks yang kolom-kolomnya merupakan basis bagi subruang,
\begin{equation*}
  A=\begin{mat}[r]
      0  &1  \\
      1  &0  \\
      0  &-1
    \end{mat}
\end{equation*}
lalu hitung.
\begin{align*}
  A\bigl(\trans{A}A\bigr)^{-1}\trans{A}
  &=\begin{mat}[r]
      0  &1  \\
      1  &0  \\
      0  &-1
    \end{mat}
    \begin{mat}[r]
      1    &0  \\
      0    &1/2
    \end{mat}
    \begin{mat}[r]
      0   &1  &0  \\
      1   &0  &-1 
    \end{mat}             \\
    &=
    \begin{mat}[r]
      1/2     &0  &-1/2  \\
      0     &1    &0   \\
      -1/2  &0    &1/2
    \end{mat}
\end{align*}
Dengan matriks tersebut, menghitung proyeksi ortogonal sebarang vektor ke $P$
menjadi mudah.
\begin{equation*}
  \proj{\vec{v}}{P}=
    \begin{mat}[r]
      1/2     &0  &-1/2  \\
      0     &1    &0   \\
      -1/2  &0    &1/2
    \end{mat}
    \colvec[r]{1 \\ -1 \\ 1}
    =\colvec[r]{0 \\ -1 \\ 0}
\end{equation*}
Sebagai pemeriksaan, perhatikan bahwa hasil ini memang berada dalam $P$.
\end{example}

\begin{exercises}
  \recommended \item
    Proyeksikan vektor-vektor ke \( M \) sepanjang \( N \).
     \begin{exparts}
       \partsitem \( \colvec[r]{3 \\ -2},\quad
                M=\set{\colvec{x \\ y}\suchthat x+y=0},\quad
                N=\set{\colvec{x \\ y}\suchthat -x-2y=0} \)
       \partsitem \( \colvec[r]{1 \\ 2},\quad
                M=\set{\colvec{x \\ y}\suchthat x-y=0},\quad
                N=\set{\colvec{x \\ y}\suchthat 2x+y=0} \)
       \partsitem \( \colvec[r]{3 \\ 0 \\ 1},\quad
                M=\set{\colvec{x \\ y \\ z}\suchthat x+y=0},\quad
                N=\set{c\cdot\colvec[r]{1 \\ 0 \\ 1}\suchthat c\in\Re} \)
     \end{exparts}
     \begin{answer}
       \begin{exparts}
         \partsitem Ketika basis-basis bagi subruang
           \begin{equation*}
             B_M=\sequence{\colvec[r]{1 \\ -1}}
             \qquad
             B_N=\sequence{\colvec[r]{2 \\ -1}}
           \end{equation*}
           dikonkatenasikan,
           \begin{equation*}
             B=\cat{B_M}{B_N}=\sequence{\colvec[r]{1 \\ -1},
                                        \colvec[r]{2 \\ -1}}
           \end{equation*}
           dan vektor yang diberikan direpresentasikan,
           \begin{equation*}
             \colvec[r]{3 \\ -2}=1\cdot\colvec[r]{1 \\ -1}+1\cdot\colvec[r]{2 \\ -1}
           \end{equation*}
           jawabannya diperoleh dengan mempertahankan bagian~$M$ dan membuang
           bagian~$N$.
           \begin{equation*}
             \proj{\colvec[r]{3 \\ -2}}{M,N}
             =\colvec[r]{1 \\ -1}
           \end{equation*}
         \partsitem Ketika basis-basis
           \begin{equation*}
             B_M=\sequence{\colvec[r]{1 \\ 1}}
             \qquad
             B_N=\sequence{\colvec[r]{1 \\ -2}}
           \end{equation*}
           dikonkatenasikan dan vektornya direpresentasikan,
           \begin{equation*}
             \colvec[r]{1 \\ 2}
               =(4/3)\cdot\colvec[r]{1 \\ 1}
                -(1/3)\cdot\colvec[r]{1 \\ -2}
           \end{equation*}
           mempertahankan hanya bagian~$M$ memberikan jawaban berikut.
           \begin{equation*}
             \proj{\colvec[r]{1 \\ 2}}{M,N}=\colvec[r]{4/3 \\ 4/3}
           \end{equation*}
         \partsitem Dengan basis-basis berikut,
           \begin{equation*}
             B_M=\sequence{\colvec[r]{1 \\ -1 \\ 0},
                           \colvec[r]{0 \\ 0 \\ 1}}
             \qquad
             B_N=\sequence{\colvec[r]{1 \\ 0 \\ 1}}
           \end{equation*}
           representasi terhadap konkatenasinya adalah berikut.
           \begin{equation*}
             \colvec[r]{3 \\ 0 \\ 1}
               =0\cdot\colvec[r]{1 \\ -1 \\ 0}
                -2\cdot\colvec[r]{0 \\ 0 \\ 1}
                +3\cdot\colvec[r]{1 \\ 0 \\ 1}
           \end{equation*}
           sehingga proyeksinya adalah berikut.
           \begin{equation*}
             \proj{\colvec[r]{3 \\ 0 \\ 1}}{M,N}=\colvec[r]{0 \\ 0 \\ -2}
           \end{equation*}
       \end{exparts}
     \end{answer}
  \recommended \item 
     Temukan \( M^\perp \).
     \begin{exparts*}
       \partsitem \( M=\set{\colvec{x \\ y}\suchthat x+y=0 } \)
       \partsitem \( M=\set{\colvec{x \\ y}\suchthat -2x+3y=0 } \)
       \partsitem \( M=\set{\colvec{x \\ y}\suchthat x-y=0 } \)
       \partsitem \( M=\set{\zero\,} \)
       \partsitem \( M=\set{\colvec{x \\ y}\suchthat x=0 } \)
       \partsitem \( M=\set{\colvec{x \\ y \\ z}\suchthat -x+3y+z=0 } \)
       \partsitem \( M=\set{\colvec{x \\ y \\ z}\suchthat 
                          \text{$x=0$ dan $y+z=0$} } \)
     \end{exparts*}
     \begin{answer}
       Seperti dalam \nearbyexample{ex:OrthoCompOne}, kita dapat menyederhanakan
       perhitungan dengan cukup mencari ruang vektor yang tegak lurus terhadap
       semua vektor dalam basis $M$.
       \begin{exparts}
         \partsitem Parameterisasi yang memberikan
           \begin{equation*}
             M=\set{c\cdot\colvec[r]{-1 \\ 1}\suchthat c\in\Re}
           \end{equation*}
           menghasilkan
           \begin{equation*}
             M^\perp=
              \set{\colvec{u \\ v}
                   \suchthat 
                      0=\colvec{u \\ v}\dotprod\colvec[r]{-1 \\ 1}}
              =\set{\colvec{u \\ v}\suchthat 0=-u+v}
           \end{equation*}
           Memparameterkan sistem linear satu-persamaan tersebut memberikan
           deskripsi berikut.
           \begin{equation*}
             M^\perp=\set{k\cdot\colvec[r]{1 \\ 1}\suchthat k\in\Re}
           \end{equation*}
         \partsitem Seperti dalam jawaban bagian sebelumnya, kita dapat
           mendeskripsikan $M$ sebagai suatu rentang,
           \begin{equation*}
             M=\set{c\cdot\colvec[r]{3/2 \\ 1}\suchthat c\in\Re}
             \qquad
             B_M=\sequence{\colvec[r]{3/2 \\ 1}}
           \end{equation*}
           lalu $M^\perp$ adalah himpunan vektor yang tegak lurus terhadap satu
           vektor dalam basis ini.
           \begin{equation*}
             M^\perp
             =\set{\colvec{u \\ v}\suchthat (3/2)\cdot u+1\cdot v=0}
             =\set{k\cdot\colvec[r]{-2/3 \\ 1}\suchthat k\in\Re}
           \end{equation*}
         \partsitem Memparameterkan syarat linear dalam deskripsi $M$
            memberikan basis berikut.
            \begin{equation*}
              M=\set{c\cdot\colvec[r]{1 \\ 1}\suchthat c\in\Re}
              \qquad
              B_M=\sequence{\colvec[r]{1 \\ 1}}
            \end{equation*}
            Sekarang, $M^\perp$ adalah himpunan vektor yang tegak lurus terhadap
            (satu-satunya vektor dalam) $B_M$.
            \begin{equation*}
              M^\perp
              =\set{\colvec{u \\ v}\suchthat u+v=0}
              =\set{k\cdot\colvec[r]{-1 \\ 1}\suchthat k\in\Re}
            \end{equation*}
            (Kebetulan, jawaban ini sesuai dengan bagian pertama soal ini.)
          \partsitem Setiap vektor dalam ruang tegak lurus terhadap vektor nol,
            sehingga $M^\perp=\Re^n$.
          \partsitem Deskripsi dan basis yang sesuai bagi $M$ bersifat rutin.
            \begin{equation*}
              M=\set{y\cdot\colvec[r]{0 \\ 1}\suchthat y\in\Re}
              \qquad
              B_M=\sequence{\colvec[r]{0 \\ 1}}
            \end{equation*}
            Kemudian
            \begin{equation*}
              M^\perp
              =\set{\colvec{u \\ v}\suchthat 0\cdot u+1\cdot v=0}
              =\set{k\cdot\colvec[r]{1 \\ 0}\suchthat k\in\Re}
            \end{equation*}
            sehingga $(\text{sumbu~$y$})^\perp=\text{sumbu~$x$}$.
          \partsitem Deskripsi $M$ mudah ditemukan melalui parameterisasi.
            \begin{equation*}
              M=\set{c\cdot\colvec[r]{3 \\ 1 \\ 0}
                     +d\cdot\colvec[r]{1 \\ 0 \\ 1}\suchthat c,d\in\Re}
              \qquad
              B_M=\sequence{\colvec[r]{3 \\ 1 \\ 0},\colvec[r]{1 \\ 0 \\ 1}}
            \end{equation*}
            Mencari $M^\perp$ di sini hanya memerlukan penyelesaian sistem
            linear dengan dua persamaan,
            \begin{equation*}
              \begin{linsys}{3}
                3u  &+  &v  &   &   &=  &0  \\
                 u  &   &   &+  &w  &=  &0
              \end{linsys}
              \grstep{-(1/3)\rho_1+\rho_2}
              \begin{linsys}{3}
                3u  &+  &v         &   &   &=  &0  \\
                    &   &-(1/3)v   &+  &w  &=  &0
              \end{linsys}
            \end{equation*}
            lalu parameterisasi.
            \begin{equation*}
              M^\perp
              =\set{k\cdot\colvec[r]{-1 \\ 3 \\ 1}\suchthat k\in\Re}
            \end{equation*}
          \partsitem Di sini, $M$ berdimensi satu,
            \begin{equation*}
              M=\set{c\cdot\colvec[r]{0 \\ -1 \\ 1}\suchthat c\in\Re}
              \qquad
              B_M=\sequence{\colvec[r]{0 \\ -1 \\ 1}}
            \end{equation*}
            sehingga $M^\perp$ berdimensi dua.
            \begin{equation*}
              M^\perp
              =\set{\colvec{u \\ v \\ w}
                     \suchthat 0\cdot u-1\cdot v+1\cdot w=0}
              =\set{j\cdot\colvec[r]{1 \\ 0 \\ 0}
                    +k\cdot\colvec[r]{0 \\ 1 \\ 1}\suchthat j,k\in\Re}
            \end{equation*}
       \end{exparts}
     \end{answer}
  \recommended \item Temukan proyeksi ortogonal vektor tersebut ke subruang.
    \begin{equation*}
      \colvec{1 \\ 2 \\ 0} 
      \qquad
      S=\spanof{\,\set{\colvec{0 \\ 2 \\ 0},\colvec{1 \\ -1 \\ 1}}\,}
    \end{equation*}
    \begin{answer}
      Dengan
      \begin{equation*}
        A=
        \begin{mat}
          0 &1 \\
          2 &-1 \\
          0 &1
        \end{mat}
      \end{equation*}
      perhitungan yang langsung, meskipun melelahkan, memberikan
      % sage: A = matrix(QQ, [[0,1], [2,-1], [0,1]])
      % sage: A*(A.transpose()*A).inverse()*A.transpose()
      % [1/2    0  1/2]
      % [  0    1    0]
      % [1/2    0  1/2]
      \begin{equation*}
        A\bigl(\trans{A}A\bigr)^{-1}\trans{A}=
        \begin{mat}
          1/2 &0 &1/2 \\
           0  &1 &0    \\
          1/2 &0 &1/2
        \end{mat}
      \end{equation*}
      Ketika diterapkan pada vektor, kita memperoleh proyeksinya.
      \begin{equation*}
        \begin{mat}
          1/2 &0 &1/2 \\
           0  &1 &0    \\
          1/2 &0 &1/2
        \end{mat}
        \colvec{1 \\ 2 \\ 0}
        =
        \colvec{1/2 \\ 2 \\  1/2}        
      \end{equation*}
      % sage: v = vector(QQ, [1, 2, 0])
      % sage: v
      % (1, 2, 0)
      % sage: A*(A.transpose()*A).inverse()*A.transpose()*v
      % (1/2, 2, 1/2)
    \end{answer}
  \recommended \item Dengan subruang yang sama seperti dalam soal sebelumnya,
    temukan proyeksi ortogonal vektor berikut.
    \begin{equation*}
      \colvec{1 \\ 2 \\ -1}
    \end{equation*}
    \begin{answer}
      Dengan menggunakan perhitungan matriks dari jawaban sebelumnya,
      \begin{equation*}
        \begin{mat}
          1/2 &0 &1/2 \\
           0  &1 &0    \\
          1/2 &0 &1/2
        \end{mat}
        \colvec{1 \\ 2 \\ -1}
        =
        \colvec{0 \\ 2 \\ 0}
      \end{equation*}
      menunjukkan proyeksi vektor tersebut.
    \end{answer}
  \recommended \item \label{exer:ProjMinimizesDistSubspace}
    Misalkan $\vec{p}$ adalah proyeksi ortogonal
    $\vec{v}\in\Re^n$ ke subruang~$S$.
    Tunjukkan bahwa $\vec{p}$ adalah titik dalam subruang yang paling dekat
    dengan~$\vec{v}$.
    \begin{answer}
      Misalkan $\vec{w}\in S$.
      Karena ini merupakan proyeksi ortogonal, kedua vektor
      $\vec{v}-\vec{p}$ dan $\vec{p}-\vec{w}$ saling tegak lurus.
      Teorema Pythagoras berlaku, sehingga sisi miring $\vec{v}-\vec{w}$
      setidaknya sepanjang $\vec{v}-\vec{p}$.
    \end{answer}
  \item 
    Subbagian ini menunjukkan dua cara melakukan proyeksi ortogonal: metode
    \nearbyexample{ex:OrthProjOne} dan~\ref{ex:ProjIntoMAlongNGener}, serta
    metode \nearbytheorem{th:OrthProjMat}.
    Untuk membandingkan keduanya, tinjau bidang $P$ yang ditentukan oleh
    $3x+2y-z=0$ dalam $\Re^3$.
    \begin{exparts}
      \partsitem Temukan basis bagi $P$.
      \partsitem Temukan $P^\perp$ dan basis bagi $P^\perp$.
      \partsitem Representasikan vektor berikut terhadap konkatenasi kedua
        basis dari bagian sebelumnya.
        \begin{equation*}
          \vec{v}=\colvec[r]{1 \\ 1 \\ 2}
        \end{equation*}
      \partsitem Temukan proyeksi ortogonal $\vec{v}$ ke $P$ dengan hanya
        mempertahankan bagian $P$ dari bagian sebelumnya.
      \partsitem Bandingkan hasil itu dengan hasil penerapan
        \nearbytheorem{th:OrthProjMat}.
    \end{exparts}
    \begin{answer}
      \begin{exparts}
        \partsitem Memparameterkan persamaan tersebut menghasilkan basis bagi
          $P$ berikut.
          \begin{equation*}
            B_P=\sequence{\colvec[r]{1 \\ 0 \\ 3},
                          \colvec[r]{0 \\ 1 \\ 2} }
          \end{equation*}
        \partsitem Karena $\Re^3$ berdimensi tiga dan $P$ berdimensi dua,
          komplemen $P^\perp$ pasti berupa garis.
          Perhitungan seperti dalam \nearbyexample{ex:OrthoCompOne},
          \begin{equation*}
            P^\perp=\set{\colvec{x \\ y \\ z}
                         \suchthat
                         \begin{mat}[r]
                           1  &0 &3  \\
                           0  &1 &2
                         \end{mat}
                         \colvec{x \\ y \\ z}
                         =\colvec[r]{0 \\ 0}\;}
          \end{equation*}
          memberikan basis bagi $P^\perp$ berikut.
          \begin{equation*}
            B_{P^\perp}=\sequence{\colvec[r]{3 \\ 2 \\ -1} }
          \end{equation*}
        \partsitem $ \displaystyle       
             \colvec[r]{1 \\ 1 \\ 2}=
             (5/14)\cdot\colvec[r]{1 \\ 0 \\ 3}+
             (8/14)\cdot\colvec[r]{0 \\ 1 \\ 2}+
             (3/14)\cdot\colvec[r]{3 \\ 2 \\ -1}$
        \partsitem $\proj{\colvec[r]{1 \\ 1 \\ 2}}{P}
          =\colvec[r]{5/14 \\ 8/14 \\ 31/14}$
        \partsitem Matriks proyeksinya,
          \begin{multline*}
            \begin{mat}[r]
              1  &0 \\
              0  &1 \\
              3  &2
            \end{mat}
            \bigl(
              \begin{mat}[r]
                1  &0  &3 \\
                0  &1  &2 
              \end{mat}
              \begin{mat}[r]
                1  &0 \\
                0  &1 \\
                3  &2
              \end{mat}
            \bigr)^{-1}
            \begin{mat}[r]
              1  &0  &3 \\
              0  &1  &2 
            \end{mat}                               \\
            \begin{aligned}
            &=
            \begin{mat}[r]
              1  &0 \\
              0  &1 \\
              3  &2
            \end{mat}
              \begin{mat}[r]
                10  &6  \\
                6   &5
              \end{mat}^{-1}
            \begin{mat}[r]
              1  &0  &3 \\
              0  &1  &2 
            \end{mat}                                     \\
            &=
            \frac{1}{14}
            \begin{mat}[r]
              5  &-6 &3  \\
             -6  &10 &2  \\
              3  &2  &13
            \end{mat}
            \end{aligned}
          \end{multline*}
          ketika diterapkan pada vektor, menghasilkan hasil yang diharapkan.
          \begin{equation*}
            \frac{1}{14}
            \begin{mat}[r]
              5  &-6 &3  \\
             -6  &10 &2  \\
              3  &2  &13
            \end{mat}
            \colvec[r]{1 \\ 1 \\ 2}
            =\colvec[r]{5/14 \\ 8/14 \\ 31/14}
          \end{equation*}
      \end{exparts}
    \end{answer}
  \item 
    Kita memiliki tiga cara untuk mencari proyeksi ortogonal suatu vektor ke
    sebuah garis: cara \nearbydefinition{def:ProjIntoLine} dari subbagian
    pertama bagian ini; cara \nearbyexample{ex:OrthProjOne}
    dan~\ref{ex:ProjIntoMAlongNGener}, yaitu merepresentasikan vektor terhadap
    suatu basis bagi ruang lalu mempertahankan bagian~$M$; dan cara
    \nearbytheorem{th:OrthProjMat}.
    Untuk kasus-kasus berikut, gunakan ketiga cara tersebut.
     \begin{exparts}
       \partsitem \( \vec{v}=\colvec[r]{1 \\ -3},\quad
                M=\set{\colvec{x \\ y}\suchthat x+y=0} \)
       \partsitem \( \vec{v}=\colvec[r]{0 \\ 1 \\ 2},\quad
                M=\set{\colvec{x \\ y \\ z}
                        \suchthat \text{$x+z=0$ dan $y=0$}} \)
     \end{exparts}
     \begin{answer}
       \begin{exparts}
         \partsitem Parameterisasi memberikan
           \begin{equation*}
             M=\set{c\cdot\colvec[r]{-1 \\ 1}\suchthat c\in\Re}
           \end{equation*}
           
           Untuk cara pertama, kita ambil vektor yang merentang garis $M$ sebagai
           \begin{equation*}
             \vec{s}=\colvec[r]{-1 \\ 1}
           \end{equation*}
           dan rumus \nearbydefinition{def:ProjIntoLine} memberikan
           \begin{equation*}
             \proj{\colvec[r]{1 \\ -3}}{\spanof{\vec{s}\,}}
              =\frac{\colvec[r]{1 \\ -3}\dotprod\colvec[r]{-1 \\ 1}}{
                     \colvec[r]{-1 \\ 1}\dotprod\colvec[r]{-1 \\ 1}}
                \cdot\colvec[r]{-1 \\ 1}
              =\frac{-4}{2}
                \cdot\colvec[r]{-1 \\ 1}
              =\colvec[r]{2 \\ -2}
           \end{equation*}

           Untuk cara kedua, kita tetapkan
           \begin{equation*}
             B_M=\sequence{\colvec[r]{-1 \\ 1}}
           \end{equation*}
           sehingga (seperti dalam \nearbyexample{ex:OrthoCompOne}
           dan~\ref{ex:OrthoCompTwo}, kita cukup mencari vektor-vektor yang
           tegak lurus terhadap semua anggota basis)
           \begin{equation*}
             M^\perp
             =\set{\colvec{u \\ v}\suchthat -1\cdot u+1\cdot v=0}
             =\set{k\cdot\colvec[r]{1 \\ 1}\suchthat k\in\Re}
             \qquad
             B_{M^\perp}=\sequence{\colvec[r]{1 \\ 1}}
           \end{equation*}
           dan merepresentasikan vektor terhadap konkatenasi tersebut
           memberikan
           \begin{equation*}
             \colvec[r]{1 \\ -3}=-2\cdot\colvec[r]{-1 \\ 1}-1\cdot\colvec[r]{1 \\ 1}
           \end{equation*}
           Mempertahankan bagian~$M$ menghasilkan jawabannya.
           \begin{equation*}
             \proj{\colvec[r]{1 \\ -3}}{M,M^\perp}=\colvec[r]{2 \\ -2}
           \end{equation*}

           Cara ketiga juga berupa perhitungan sederhana
           (terdapat matriks $\nbyn{1}$ di tengah, dan inversnya juga
           berukuran $\nbyn{1}$),
           \begin{multline*}
             A\left(\trans{A}A\right)^{-1}\trans{A}           \\
             \begin{aligned}
             &=
             \begin{mat}[r]
               -1 \\ 1
             \end{mat}
             \left(
               \begin{mat}[r]
                 -1  &1
               \end{mat}
               \begin{mat}[r]
                 -1  \\
                  1
               \end{mat}
             \right)^{-1}
             \begin{mat}[r]
               -1  &1
             \end{mat}
             =
             \begin{mat}[r]
               -1 \\ 1
             \end{mat}
               \begin{mat}[r]
                 2
               \end{mat}^{-1}
             \begin{mat}[r]
               -1  &1
             \end{mat}               \\
             &=
             \begin{mat}[r]
               -1 \\ 1
             \end{mat}
               \begin{mat}[r]
                 1/2
               \end{mat}
             \begin{mat}[r]
               -1  &1
             \end{mat}               
             =
             \begin{mat}[r]
               -1 \\ 1
             \end{mat}
             \begin{mat}[r]
               -1/2  &1/2
             \end{mat}               
             =
             \begin{mat}[r]
               1/2  &-1/2  \\
              -1/2  &1/2  
             \end{mat}
             \end{aligned}
           \end{multline*}
           yang tentu saja memberikan jawaban yang sama.
           \begin{equation*}
             \proj{\colvec[r]{1 \\ -3}}{M}
             =
             \begin{mat}[r]
               1/2  &-1/2  \\
              -1/2  &1/2  
             \end{mat}
             \colvec[r]{1 \\ -3}
             =\colvec[r]{2 \\ -2}
           \end{equation*}
         \partsitem Parameterisasi memberikan
           \begin{equation*}
             M=\set{c\cdot\colvec[r]{-1 \\ 0 \\ 1}\suchthat c\in\Re}
           \end{equation*}
           Dengan itu, rumus bagi cara pertama memberikan
           \begin{equation*}
             \frac{\colvec[r]{0 \\ 1 \\ 2}\dotprod\colvec[r]{-1 \\ 0 \\ 1}}{
                   \colvec[r]{-1 \\ 0 \\ 1}\dotprod\colvec[r]{-1 \\ 0 \\ 1}}
               \cdot\colvec[r]{-1 \\ 0 \\ 1}
             =\frac{2}{2}
               \cdot\colvec[r]{-1 \\ 0 \\ 1}
             =\colvec[r]{-1 \\ 0 \\ 1}
           \end{equation*}
           Untuk melanjutkan dengan metode kedua, kita mencari $M^\perp$,
           \begin{equation*}
             M^\perp
             =\set{\colvec{u \\ v \\ w}\suchthat -u+w=0}
             =\set{j\cdot\colvec[r]{1 \\ 0 \\ 1}
                   +k\cdot\colvec[r]{0 \\ 1 \\ 0}\suchthat j,k\in\Re}
           \end{equation*}
           mencari representasi vektor yang diberikan terhadap konkatenasi
           basis $B_M$ dan $B_{M^\perp}$,
           \begin{equation*}
             \colvec[r]{0 \\ 1 \\ 2}
              =1\cdot\colvec[r]{-1 \\ 0 \\ 1}
               +1\cdot\colvec[r]{1 \\ 0 \\ 1}
               +1\cdot\colvec[r]{0 \\ 1 \\ 0}
           \end{equation*}
           dan hanya mempertahankan bagian~$M$.
           \begin{equation*}
             \proj{\colvec[r]{0 \\ 1 \\ 2}}{M}=1\cdot\colvec[r]{-1 \\ 0 \\ 1}
                       =\colvec[r]{-1 \\ 0 \\ 1}
           \end{equation*}
           Akhirnya, untuk metode ketiga, perhitungan matriks
           \begin{multline*}
             A\left(\trans{A}A\right)^{-1}\trans{A}                 \\
             \begin{aligned}
             &=
             \begin{mat}[r]
               -1 \\ 0  \\  1
             \end{mat}
             \bigl(
               \begin{mat}[r]
                 -1  &0  &1
               \end{mat}
               \begin{mat}[r]
                 -1  \\
                  0  \\
                  1
               \end{mat}
             \bigr)^{-1}
             \begin{mat}[r]
               -1  &0  &1
             \end{mat}
             =
             \begin{mat}[r]
               -1 \\ 0  \\  1
             \end{mat}
               \begin{mat}[r]
                 2
               \end{mat}^{-1}
             \begin{mat}[r]
               -1  &0  &1
             \end{mat}              \\
             &=
             \begin{mat}[r]
               -1 \\ 0  \\  1
             \end{mat}
               \begin{mat}[r]
                 1/2
               \end{mat}
             \begin{mat}[r]
               -1  &0  &1
             \end{mat}
             =
             \begin{mat}[r]
               -1 \\ 0  \\  1
             \end{mat}
             \begin{mat}[r]
               -1/2  &0  &1/2
             \end{mat}                                 \\
             &=
             \begin{mat}[r]
               1/2  &0  &-1/2  \\
               0    &0  &0     \\
              -1/2  &0  &1/2
             \end{mat}
             \end{aligned}
           \end{multline*}
           yang diikuti oleh perkalian matriks-vektor
           \begin{equation*}
             \proj{\colvec[r]{0 \\ 1 \\ 2}}{M}=
             \begin{mat}[r]
               1/2  &0  &-1/2  \\
               0    &0  &0     \\
              -1/2  &0  &1/2
             \end{mat}
             \colvec[r]{0 \\ 1 \\ 2}
             =\colvec[r]{-1 \\ 0 \\ 1}
           \end{equation*}
           memberikan jawabannya.
       \end{exparts}
     \end{answer}
  \item 
    Periksa bahwa operasi dalam \nearbydefinition{def:ProjIntoMAlongN}
    terdefinisi dengan baik.
    Artinya, dalam \nearbyexample{ex:OrthProjOne}
    dan~\ref{ex:ProjIntoMAlongNGener}, bukankah jawabannya bergantung pada
    pilihan basis?
    \begin{answer}
      Tidak. Penguraian vektor $\vec{v}=\vec{m}+\vec{n}$ menjadi
      $\vec{m}\in M$ dan $\vec{n}\in N$ tidak bergantung pada basis yang dipilih
      bagi subruang tersebut, seperti yang kita tunjukkan dalam subbagian
      Jumlah Langsung.
    \end{answer}
  \item 
    Apa hasil proyeksi ortogonal ke subruang trivial?
    \begin{answer}
      Proyeksi ortogonal suatu vektor ke sebuah subruang merupakan anggota
      subruang tersebut.
      Karena subruang trivial hanya memiliki satu anggota, $\zero$, proyeksi
      sebarang vektor harus sama dengan $\zero$.
    \end{answer}
  \item 
    Apa hasil proyeksi \( \vec{v} \) ke \( M \) sepanjang \( N \) jika
    \( \vec{v}\in M \)?
    \begin{answer}
      Proyeksi suatu $\vec{v}\in M$ ke $M$ sepanjang $N$ adalah $\vec{v}$.
      Menguraikan $\vec{v}=\vec{m}+\vec{n}$ memberikan $\vec{m}=\vec{v}$ dan
      $\vec{n}=\zero$; membuang bagian~$N$ tetapi mempertahankan bagian~$M$
      menghasilkan proyeksi $\vec{m}=\vec{v}$.
    \end{answer}
  \item
    Tunjukkan bahwa jika \( M\subseteq\Re^n \) merupakan subruang dengan basis
    ortonormal \( \basis{\kappa}{k} \), maka proyeksi ortogonal
    \( \vec{v} \) ke \( M \) adalah berikut.
    \begin{equation*}
      (\vec{v}\dotprod\vec{\kappa}_1)\cdot\vec{\kappa}_1+
      \dots+
      (\vec{v}\dotprod\vec{\kappa}_k)\cdot\vec{\kappa}_k
    \end{equation*}
    \begin{answer}
      Bukti \nearbylemma{le:OrthoProjWellDefd} menunjukkan bahwa setiap vektor
      $\vec{v}\in\Re^n$ merupakan jumlah proyeksi ortogonalnya ke garis-garis
      yang direntang oleh vektor-vektor basis.
      \begin{equation*}
        \vec{v}=\proj{\vec{v}\,}{\spanof{\vec{\kappa}_1}}
                 +\dots+\proj{\vec{v}\,}{\spanof{\vec{\kappa}_k}}
               =\frac{\vec{v}\dotprod\vec{\kappa}_1}{
                      \vec{\kappa}_1\dotprod\vec{\kappa}_1}
                   \cdot\vec{\kappa}_1
               +\dots
               +\frac{\vec{v}\dotprod\vec{\kappa}_k}{
                      \vec{\kappa}_k\dotprod\vec{\kappa}_k}
                   \cdot\vec{\kappa}_k
      \end{equation*}
      Karena basisnya ortonormal, penyebut setiap pecahan memenuhi
      $\vec{\kappa}_i\dotprod\vec{\kappa}_i=1$.
    \end{answer}
  \recommended \item
    Buktikan bahwa pemetaan \( \map{p}{V}{V} \) merupakan proyeksi ke \( M \)
    sepanjang \( N \) jika dan hanya jika pemetaan \( \identity-p \) merupakan
    proyeksi ke \( N \) sepanjang \( M \).
    (Ingat definisi selisih dua pemetaan:
    $(\identity-p)\,(\vec{v})=\identity(\vec{v})-p(\vec{v})
     =\vec{v}-p(\vec{v})$.)
     \begin{answer}
       Jika $V=M\directsum N$, maka setiap vektor terurai secara tunggal sebagai
       $\vec{v}=\vec{m}+\vec{n}$.
       Bagi semua $\vec{v}$, pemetaan $p$ memberikan
       $p(\vec{v})=\vec{m}$ jika dan hanya jika
       $\vec{v}-p(\vec{v})=\vec{n}$, sebagaimana yang disyaratkan.
     \end{answer}
  \item \label{exer:PerpOnBasisImplPerp} 
    Tunjukkan bahwa jika suatu vektor tegak lurus terhadap setiap vektor dalam
    suatu himpunan, maka vektor itu tegak lurus terhadap setiap vektor dalam
    rentang himpunan tersebut.
    \begin{answer}
      Misalkan $\vec{v}$ tegak lurus terhadap setiap $\vec{w}\in S$.
      Maka
      $\vec{v}\dotprod(c_1\vec{w}_1+\dots+c_n\vec{w}_n)
       =\vec{v}\dotprod(c_1\vec{w}_1)
          +\dots+\vec{v}\dotprod (c_n\vec{w}_n)
       =c_1(\vec{v}\dotprod\vec{w}_1)
          +\dots+c_n(\vec{v}\dotprod\vec{w}_n)
       =c_1\cdot 0+\dots+c_n\cdot 0=0$.
    \end{answer}
  \item
    Benar atau salah:~irisan suatu subruang dan komplemen ortogonalnya bersifat
    trivial.
    \begin{answer}
      Benar; satu-satunya vektor yang ortogonal terhadap dirinya sendiri adalah
      vektor nol.
    \end{answer}
  \item 
    Tunjukkan bahwa jumlah dimensi komplemen-komplemen ortogonal sama dengan
    dimensi seluruh ruang.
    \begin{answer}
      Ini langsung mengikuti pernyataan dalam
      \nearbylemma{le:OrthoProjWellDefd} bahwa ruang merupakan jumlah langsung
      keduanya.
    \end{answer}
  \item
    Misalkan \( \vec{v}_1,\vec{v}_2\in\Re^n \) sedemikian sehingga bagi semua
    subruang komplementer \( M,N\subseteq\Re^n \), proyeksi \( \vec{v}_1 \)
    dan \( \vec{v}_2 \) ke \( M \) sepanjang \( N \) sama.
    Haruskah \( \vec{v}_1 \) sama dengan \( \vec{v}_2 \)?
    (Jika ya, bagaimana jika syaratnya diperlonggar menjadi:~semua proyeksi
     ortogonal keduanya sama?)
    \begin{answer}
      Keduanya harus sama, bahkan hanya di bawah syarat yang tampaknya lebih
      lemah bahwa keduanya menghasilkan hasil yang sama pada semua proyeksi
      ortogonal.
      Tinjau subruang $M$ yang direntang oleh himpunan
      $\set{\vec{v}_1,\vec{v}_2}$.
      Karena keduanya berada dalam $M$, proyeksi ortogonal $\vec{v}_1$ ke $M$
      adalah $\vec{v}_1$, dan proyeksi ortogonal $\vec{v}_2$ ke $M$ adalah
      $\vec{v}_2$.
      Agar proyeksi keduanya ke $M$ sama, keduanya harus sama.
    \end{answer}
  \recommended \item \label{exer:AlgOfPerps}
    Misalkan \( M,N \) merupakan subruang-subruang \( \Re^n \).
    Operator perp bertindak pada subruang; kita dapat menanyakan bagaimana
    operator ini berinteraksi dengan operasi-operasi lain semacam itu.
    \begin{exparts}
      \partsitem Tunjukkan bahwa dua perp saling meniadakan:
        \( (M^\perp)^\perp=M \).
      \partsitem %Schaums p331 #114
        Buktikan bahwa \( M\subseteq N \) mengakibatkan
           bahwa \( N^\perp\subseteq M^\perp \).
      \partsitem Tunjukkan bahwa
        \( (M+N)^\perp=M^\perp\intersection N^\perp \).
      %\partsitem Show that \( (M\intersection N)^\perp=M^\perp +N^\perp \).
    \end{exparts}
    \begin{answer}
      \begin{exparts}
        \partsitem 
          Kita akan menunjukkan bahwa kedua himpunan saling memuat,
          $M\subseteq (M^\perp)^\perp$ dan $(M^\perp)^\perp \subseteq M$.
          Untuk inklusi pertama, jika $\vec{m}\in M$, maka berdasarkan definisi
          operasi perp, $\vec{m}$ tegak lurus terhadap setiap
          $\vec{v}\in M^\perp$, dan oleh karena itu (sekali lagi berdasarkan
          definisi operasi perp)
          $\vec{m}\in (M^\perp)^\perp$.
          Untuk arah lainnya, tinjau $\vec{v}\in (M^\perp)^\perp$.
          Bukti \nearbylemma{le:OrthoProjWellDefd} menunjukkan bahwa
          $\Re^n=M\directsum M^\perp$ dan bahwa
          kita dapat memberikan basis ortogonal bagi ruang tersebut,
          $\sequence{\vec{\kappa}_1,\dots,\vec{\kappa}_k,
                        \vec{\kappa}_{k+1},\dots,\vec{\kappa}_n}$
          sedemikian sehingga bagian pertama
          $\sequence{\vec{\kappa}_1,\dots,\vec{\kappa}_k}$ merupakan basis
          bagi $M$, dan bagian kedua merupakan basis bagi $M^\perp$.
          Bukti tersebut juga memeriksa bahwa setiap vektor dalam ruang
          merupakan jumlah proyeksi ortogonalnya ke garis-garis yang direntang
          oleh vektor-vektor basis ini.
          \begin{equation*}
             \vec{v}=\proj{\vec{v}\,}{\spanof{\vec{\kappa}_1}}
                          +\dots+\proj{\vec{v}\,}{\spanof{\vec{\kappa}_n}}
          \end{equation*}
          Karena $\vec{v}\in (M^\perp)^\perp$, vektor itu tegak lurus terhadap
          setiap vektor dalam $M^\perp$, sehingga semua proyeksi pada bagian
          kedua bernilai nol.
          Jadi, $\vec{v}=\proj{\vec{v}\,}{\spanof{\vec{\kappa}_1}}
                          +\dots+\proj{\vec{v}\,}{\spanof{\vec{\kappa}_k}}$,
          yang merupakan kombinasi linear vektor-vektor dari $M$, sehingga
          $\vec{v}\in M$.
          (\textit{Catatan}.
           Berikut cara yang lebih ringkas untuk mengerjakan inklusi kedua:
           tuliskan ruang sebagai $M\directsum M^\perp$ dan juga sebagai
           $M^\perp\directsum (M^\perp)^\perp$.
           Karena bagian pertama menunjukkan bahwa
           $M\subseteq (M^\perp)^\perp$ dan kalimat sebelumnya menunjukkan
           bahwa kedua subruang $M$ dan $(M^\perp)^\perp$ berdimensi sama,
           dapat disimpulkan bahwa $M$ sama dengan $(M^\perp)^\perp$.)
        \partsitem Karena $M\subseteq N$, setiap $\vec{v}$ yang tegak lurus
           terhadap setiap vektor dalam $N$ juga tegak lurus terhadap setiap
           vektor dalam $M$.
           Namun, kalimat itu tepat menyatakan bahwa
           $N^\perp\subseteq M^\perp$.
        \partsitem Sekali lagi kita akan menunjukkan kesamaan kedua himpunan
           melalui inklusi dua arah.
           Arah pertama mudah; setiap $\vec{v}$ yang tegak lurus terhadap
           setiap vektor dalam
           $M+N=\set{\vec{m}+\vec{n}\suchthat \vec{m}\in M,\, \vec{n}\in N}$
           tegak lurus terhadap setiap vektor berbentuk $\vec{m}+\zero$
           (yakni, setiap vektor dalam $M$) dan setiap vektor berbentuk
           $\zero+\vec{n}$ (setiap vektor dalam $N$), sehingga
           $(M+N)^\perp\subseteq M^\perp\intersection N^\perp$. 
           Arah kedua juga rutin; sebarang vektor
           $\vec{v}\in M^\perp\intersection N^\perp$ 
           tegak lurus terhadap sebarang vektor berbentuk
           $c\vec{m}+d\vec{n}$ karena
           $\vec{v}\dotprod (c\vec{m}+d\vec{n})
             =c\cdot(\vec{v}\dotprod\vec{m})+d\cdot(\vec{v}\dotprod\vec{n})
             =c\cdot 0+d\cdot 0
             =0$.
      \end{exparts}
    \end{answer}
  \recommended \item 
    Materi dalam subbagian ini memungkinkan kita menyatakan suatu hubungan
    geometris antara ruang hasil dan ruang nol suatu pemetaan linear yang belum
    pernah kita lihat.
    \begin{exparts}
      \partsitem Representasikan $\map{f}{\Re^3}{\Re}$ yang diberikan oleh
        \begin{equation*}
          \colvec{v_1 \\ v_2 \\ v_3}
          \mapsto
          1v_1+2v_2+ 3v_3
        \end{equation*}
        terhadap basis-basis standar, lalu tunjukkan bahwa
        \begin{equation*}
          \colvec[r]{1 \\ 2 \\ 3}
        \end{equation*}
        merupakan anggota perp dari ruang nol.
        Buktikan bahwa $\nullspace{f}^\perp$ sama dengan rentang vektor ini.
      \partsitem Perumum hasil tersebut agar berlaku bagi sebarang
        $\map{f}{\Re^n}{\Re}$.
      \partsitem Representasikan $\map{f}{\Re^3}{\Re^2}$,
        \begin{equation*}
          \colvec{v_1 \\ v_2 \\ v_3}
          \mapsto
          \colvec{1v_1+2v_2+ 3v_3 \\
                  4v_1+5v_2+ 6v_3}   
        \end{equation*}
        terhadap basis-basis standar, lalu tunjukkan bahwa
        \begin{equation*}
          \colvec[r]{1 \\ 2 \\ 3},
          \;\colvec[r]{4 \\ 5 \\ 6}
        \end{equation*}
        keduanya merupakan anggota perp dari ruang nol.
        Buktikan bahwa $\nullspace{f}^\perp$ merupakan rentang keduanya.
        (\textit{Petunjuk}.
         Lihat bagian ketiga \nearbyexercise{exer:AlgOfPerps}.)
      \partsitem Perumum hasil tersebut agar berlaku bagi sebarang
        $\map{f}{\Re^n}{\Re^m}$.
    \end{exparts}
    Dalam \cite{Strang93}
    hasil ini disebut
    \definend{Teorema Fundamental Aljabar Linear}%
    \index{Teorema Fundamental!Aljabar Linear}
    \begin{answer}
      \begin{exparts}
        \partsitem Representasi
          \begin{equation*}
            \colvec{v_1 \\ v_2 \\ v_3}
            \mapsunder{f}
            1v_1+2v_2+ 3v_3    
          \end{equation*}
          adalah berikut.
          \begin{equation*}
            \rep{f}{\stdbasis_3,\stdbasis_1}
            =\begin{mat}[r]
               1  &2  &3
             \end{mat}
          \end{equation*}
          Berdasarkan definisi $f$,
          \begin{equation*}
            \nullspace{f}
            =\set{\colvec{v_1 \\ v_2 \\ v_3}
                  \suchthat 1v_1+2v_2+3v_3=0}
            =\set{\colvec{v_1 \\ v_2 \\ v_3}
                  \suchthat 
                   \colvec[r]{1 \\ 2 \\ 3}\dotprod\colvec{v_1 \\ v_2 \\ v_3}=0}
          \end{equation*}
          dan deskripsi kedua ini tepat menyatakan
          \begin{equation*}
            \nullspace{f}^\perp
            =\spanof{\set{\colvec[r]{1 \\ 2 \\ 3}}}
          \end{equation*}
        \partsitem Perumumannya ialah bahwa bagi sebarang
           $\map{f}{\Re^n}{\Re}$ terdapat vektor $\vec{h}$ sedemikian sehingga
           \begin{equation*}
             \colvec{v_1 \\ \vdots \\ v_n}
             \mapsunder{f}
             h_1v_1+\dots+h_nv_n
           \end{equation*}
           dan
           $\nullspace{f}^\perp=\spanof{\set{\vec{h}}}$.
           Seperti dalam bagian sebelumnya, kita dapat membuktikannya dengan
           merepresentasikan $f$ terhadap basis-basis standar dan mengambil
           $\vec{h}$ sebagai vektor kolom yang diperoleh dengan mentransposkan
           satu baris dari representasi matriks tersebut.
         \partsitem Tentu saja,
           \begin{equation*}
             \rep{f}{\stdbasis_3,\stdbasis_2}
             =
             \begin{mat}[r]
               1  &2  &3  \\
               4  &5  &6
             \end{mat}
           \end{equation*}
           sehingga ruang nolnya adalah himpunan berikut.
           \begin{equation*}
             \nullspace{f}=
             \set{\colvec{v_1 \\ v_2 \\ v_3}
                  \suchthat
                   \;\begin{mat}[r]
                     1  &2  &3  \\
                     4  &5  &6
                   \end{mat}
                   \colvec{v_1 \\ v_2 \\ v_3}
                   =\colvec[r]{0  \\ 0}\;}
           \end{equation*}
           Deskripsi tersebut memperjelas bahwa
           \begin{equation*}
             \colvec[r]{1 \\ 2 \\ 3},
              \,\colvec[r]{4 \\ 5 \\ 6}\in\nullspace{f}^\perp
           \end{equation*}
           dan karena $\nullspace{f}^\perp$ merupakan subruang $\Re^n$,
           rentang kedua vektor tersebut merupakan subruang dari perp ruang
           nol.
           Untuk melihat bahwa inklusi ini merupakan kesamaan, ambil
           \begin{equation*}
             M=\spanof{\set{\colvec[r]{1 \\ 2 \\ 3}}}
             \qquad
             N=\spanof{\set{\colvec[r]{4 \\ 5 \\ 6}}}
           \end{equation*}
           dalam bagian ketiga \nearbyexercise{exer:AlgOfPerps}, seperti yang
           disarankan dalam petunjuk.
        \partsitem Seperti di atas, memperumum dari kasus khusus ini mudah:~bagi
           sebarang $\map{f}{\Re^n}{\Re^m}$, matriks $H$ yang merepresentasikan
           pemetaan terhadap basis-basis standar mendeskripsikan tindakan
           \begin{equation*}
             \colvec{v_1 \\ \vdots \\ v_n}
             \mapsunder{f}
             \colvec{h_{1,1}v_1+h_{1,2}v_2+\dots+h_{1,n}v_n \\
                     \vdots                                 \\
                     h_{m,1}v_1+h_{m,2}v_2+\dots+h_{m,n}v_n}
           \end{equation*}
           dan deskripsi ruang nol memberikan bahwa dengan mentransposkan
           $m$~baris $H$,
           \begin{equation*}
             \vec{h}_1=\colvec{h_{1,1} \\ h_{1,2} \\ \vdots \\ h_{1,n}},
              \dots
             \vec{h}_m=\colvec{h_{m,1} \\ h_{m,2} \\ \vdots \\ h_{m,n}}
           \end{equation*}
           diperoleh $\nullspace{f}^\perp
                    =\spanof{\set{\vec{h}_1,\dots,\vec{h}_m}}$.
           (\cite{Strang93} mendeskripsikan ruang ini sebagai transpos ruang
           baris $H$.)
      \end{exparts}
    \end{answer}
  \item
    Definisikan suatu \definend{proyeksi\/}\index{proyeksi} sebagai transformasi
    linear \( \map{t}{V}{V} \) dengan sifat bahwa mengulangi proyeksi tidak
    menghasilkan apa pun selain yang dihasilkan oleh satu kali proyeksi:
    \( (\composed{t}{t})\,(\vec{v})=t(\vec{v}) \) bagi semua
    \( \vec{v}\in V \).
    \begin{exparts}
      \partsitem Tunjukkan bahwa proyeksi ortogonal ke sebuah garis memiliki
        sifat tersebut.
      \partsitem Tunjukkan bahwa proyeksi sepanjang suatu subruang memiliki
        sifat tersebut.
      \partsitem Tunjukkan bahwa bagi setiap $t$ semacam itu, terdapat basis
        \( B=\basis{\beta}{n} \) bagi \( V \) sedemikian sehingga
        \begin{equation*}
          t(\vec{\beta}_i)=
          \begin{cases}
                  \vec{\beta}_i  &i=1,2,\dots,\,r  \\
                  \zero          &i=r+1,r+2,\dots,\,n
                \end{cases}
        \end{equation*}
        di mana \( r \) adalah rank \( t \).
      \partsitem Simpulkan bahwa setiap proyeksi merupakan proyeksi sepanjang
        suatu subruang.
      \partsitem Simpulkan pula bahwa setiap proyeksi memiliki representasi
        \begin{equation*}
          \rep{t}{B,B}=
          \begin{pmat}{c|c}
            I   &Z  \\ \hline
            Z   &Z
          \end{pmat}
        \end{equation*}
        dalam bentuk identitas-parsial blok.
    \end{exparts}
    \begin{answer}
      \begin{exparts}
        \partsitem Pertama, perhatikan bahwa jika vektor $\vec{v}$ sudah berada
          pada garis, maka proyeksi ortogonal menghasilkan $\vec{v}$ sendiri.
          Salah satu cara memverifikasinya ialah menerapkan rumus proyeksi ke
          garis yang direntang oleh vektor $\vec{s}$, yaitu
          $(\vec{v}\dotprod\vec{s}/\vec{s}\dotprod\vec{s})\cdot\vec{s}$.
          Dengan mengambil garis sebagai
          $\set{k\cdot\vec{v}\suchthat k\in\Re}$
          (kasus $\vec{v}=\zero$ terpisah tetapi mudah), diperoleh
          $(\vec{v}\dotprod\vec{v}/\vec{v}\dotprod\vec{v})\cdot\vec{v}$,
          yang menyederhana menjadi $\vec{v}$, sebagaimana yang disyaratkan.
 
          Hal itu menjawab soal karena setelah satu kali proyeksi ke garis,
          hasil $\proj{\vec{v}}{\ell}$ berada pada garis tersebut.
          Paragraf sebelumnya menyatakan bahwa memproyeksikannya lagi ke garis
          yang sama tidak akan memberikan efek apa pun.
        \partsitem Argumen di sini serupa dengan argumen pada bagian sebelumnya.
          Dengan $V=M\directsum N$, proyeksi $\vec{v}=\vec{m}+\vec{n}$ adalah
          $\proj{\vec{v}\,}{M,N}=\vec{m}$.
          Mengulangi proyeksi kini memberikan
          $\proj{\vec{m}}{M,N}=\vec{m}$, sebagaimana yang disyaratkan, karena
          penguraian anggota $M$ menjadi jumlah anggota $M$ dan anggota $N$
          adalah $\vec{m}=\vec{m}+\zero$.
          Jadi, memproyeksikan dua kali ke $M$ sepanjang $N$ memiliki efek yang
          sama dengan memproyeksikan sekali.
        \partsitem Seperti yang ditunjukkan bagian-bagian sebelumnya, syarat
          tersebut membuat $t$ tidak mengubah vektor-vektor dalam ruang hasil,
          dan menunjukkan bahwa kita sebaiknya mengambil
          $\vec{\beta}_1$, \ldots, $\vec{\beta}_r$ sebagai vektor-vektor basis
          bagi ruang hasil; artinya, kita mengambil ruang hasil $t$ sebagai
          $M$ (sehingga $\dim(M)=r$).
          Untuk komplemennya, tuliskan $N$ sebagai ruang nol $t$, dan kita akan
          menunjukkan bahwa $V=M\directsum N$.

          Untuk menunjukkannya, kita dapat menunjukkan bahwa irisan keduanya
          trivial, $M\intersection N=\set{\zero}$, dan jumlah keduanya adalah
          seluruh ruang, $M+N=V$.
          Untuk yang pertama, jika vektor $\vec{m}$ berada dalam ruang hasil,
          maka terdapat $\vec{v}\in V$ dengan $t(\vec{v})=\vec{m}$, dan syarat
          pada $t$ memberikan
          $t(\vec{m})=(\composed{t}{t})\,(\vec{v})=t(\vec{v})=\vec{m}$.
          Jika vektor yang sama juga berada dalam ruang nol, maka
          $t(\vec{m})=\zero$, sehingga irisan ruang hasil dan ruang nol trivial.
          Untuk yang kedua, guna menuliskan sebarang $\vec{v}$ sebagai jumlah
          suatu vektor dari ruang hasil dan suatu vektor dari ruang nol, fakta
          bahwa syarat $t(\vec{v})=t(t(\vec{v}))$ dapat ditulis ulang sebagai
          $t(\vec{v}-t(\vec{v}))=\zero$ menunjukkan agar kita mengambil
          $\vec{v}=t(\vec{v})+(\vec{v}-t(\vec{v}))$.
  
          Untuk menyelesaikannya, ambil basis
          $B=\sequence{\vec{\beta}_1,\ldots,\vec{\beta}_n}$ bagi $V$, dengan
          $\sequence{\vec{\beta}_1,\ldots,\vec{\beta}_r}$ sebagai basis bagi
          ruang hasil $M$, dan
          $\sequence{\vec{\beta}_{r+1},\ldots,\vec{\beta}_n}$ sebagai basis
          bagi ruang nol $N$.
        \partsitem Setiap proyeksi (sebagaimana didefinisikan dalam latihan
          ini) merupakan proyeksi ke ruang hasilnya sepanjang ruang nolnya.
        \partsitem Ini juga langsung mengikuti bagian ketiga.
      \end{exparts}
    \end{answer}
%  \item 
%    Show that projection into a plane followed by projection from the plane
%    to a line in the plane equals the single operation of projection into the
%    line.
%  \item 
%    Assume that \( V \) decomposes as \( M_1\directsum N_1 \) and also as
%    \( M_2\directsum N_2 \).
%    Let the function \( p_1 \) 
%    be the  projection into \( M_1 \) along \( N_1 \) and let
%    \( p_2 \) be the projection into \( M_2 \) along \( N_2 \).
%    Prove that if \( \composed{p_1}{p_2} \) and \( \composed{p_2}{p_1} \) are
%    both the zero map then \( p_1+p_2 \) projects into \( M_1+M_2 \) along
%    \( N_1\intersection N_2 \).
%  \item 
%    Why does the definition of projection into \( M \) along \( N \)
%    require that the subspaces be complements?
%    Can we relax that requirement and just have
%    \( N \intersection M=\set{\zero} \) (but not have the span of
%    \( N\union M \) be the whole space)?
%    Describe the projection of \( \vec{v} \) into \( M \) along \( N \)
%    if \( M=N \).
%    What if \( M\intersection N \) is nontrivial?
%    What if it is trivial?
%  \item 
%    \begin{exparts}
%      \partsitem Show that if the columns of a matrix \( A \) form a linearly
%        independent set then \( \trans{A}A \) is invertible.
%      \partsitem Show that if the columns of \( A \) are mutually 
%        orthogonal and
%        nonzero then \( \trans{A}A \) is the identity matrix.
%    \end{exparts}
  \item 
    Suatu matriks persegi disebut \definend{simetris\/}%\index{symmetric!matrix}%
    \index{matriks!simetris} jika setiap entri \( i,j \) sama dengan entri
    \( j,i \) (yakni, jika matriks sama dengan transposnya).
    Tunjukkan bahwa matriks proyeksi \( A(\trans{A}A)^{-1}\trans{A} \)
    simetris.
    \cite{Strang}
    \textit{Petunjuk}. Temukan sifat-sifat transpos dengan melihat indeks di
      bawah `transpos'.
    \begin{answer}
      Bagi sebarang matriks $M$ berlaku
      $\trans{(M^{-1})}=(\trans{M})^{-1}$, dan bagi sebarang dua matriks $M$,
      $N$ berlaku $\trans{MN}=\trans{N}\trans{M}$ (tentu saja asalkan invers
      dan hasil kalinya terdefinisi).
      Menerapkan kedua sifat ini memberikan bahwa matriks tersebut sama dengan
      transposnya.
      \begin{multline*}
         \trans{\bigl(A(\trans{A}A)^{-1}\trans{A}\bigr)}
         =(\trans{\trans{A}})
           (\trans{\bigl((\trans{A}A)^{-1}\bigr)})
           (\trans{A})                               \\
         =(\trans{\trans{A}})
           ( \bigl(\trans{(\trans{A}A)}\bigr)^{-1} )
           (\trans{A})                    
         =A
           (\trans{A}\trans{\trans{A}})^{-1}
           \trans{A}
         =A
           (\trans{A}A)^{-1}
           \trans{A}
      \end{multline*}
    \end{answer}
\end{exercises}
